\documentclass[output=paper,colorlinks,citecolor=brown
% ,hidelinks
% showindex
]{langscibook}
\ChapterDOI{10.5281/zenodo.20285194}
\author{Rebecca Grollemund\affiliation{University of Missouri} and Gérard Philippson\affiliation{Laboratoire Dynamique Du Langage, CNRS} and Braden Brown\affiliation{University of Missouri} and  Chaeli Rule\affiliation{The Ohio State University}}

%Insert your title here
\title{A reconstruction of Proto-Sawabantu}  
\abstract{The present paper aims at studying and comparing Bantu languages spoken along the Cameroonian, Equatorial Guinean and Gabonese coasts called the Sawabantu languages by reconstructing proto-Sawabantu (consonants and vowels) via the Comparative Method used in Historical Linguistics. The proto-Sawabantu reconstructions obtained are analyzed and compared with other proto-languages reconstructed for languages spoken in this region: proto-Bantu and proto-Manenguba in order to identify Sawabantu lexical innovations and phonological innovations. The resulting historical analyses shed new light on the fascinating history of the speakers of this particular region where the great Bantu Expansion started 5,000 years ago in southern Cameroon.}


\IfFileExists{../localcommands.tex}{
   \addbibresource{../localbibliography.bib}
   %%%%%%%%%%%%%%%%%%%%%%%%%%%%%%%%%%%%%%%%%%%%%
%%%%%%%%%% From LSP %%%%%%%%%%%%%%%%%%%%%%%%%
%%%%%%%%%%%%%%%%%%%%%%%%%%%%%%%%%%%%%%%%%%%%%

\usepackage{tabularx,multicol}
\usepackage{url}
\urlstyle{same}

\usepackage{langsci-optional}
\usepackage{langsci-lgr}
\usepackage{langsci-gb4e}

% NOTE THAT THE FOLLOWING PACKAGES ARE BANNED: 
% - pstricks
% - tipa
%%%%%%%%%%%%%%%%%%%%%%%%%%%%%%%%%%%%%%%%%%%%%
%%%%%%%%%% From ACAL LaTeX template %%%%%%%%%
%%%%%%%%%%%%%%%%%%%%%%%%%%%%%%%%%%%%%%%%%%%%%

%For stacking text, used here in autosegmental diagrams
\usepackage{stackengine}

%To combine rows in tables
\usepackage{multirow}

%Gives some extra formatting options, e.g. underlining/strikeout
\usepackage[normalem]{ulem}

%Use package/command below to create a double-spaced document, if you want one. Uncomment BOTH the package and the command (\doublespacing) to create a doublespaced document, or leave them as is to have a single-spaced document.
%\usepackage{setspace}
%\doublespacing 

 

%use for special OT tableaux symbols like bomb and sad face. must be loaded early on because it doesn't play well with some other packages
\usepackage{fourier-orns}

%Basic math symbols 
\usepackage{pifont}
\usepackage{amssymb}

%%%Gives shortcuts for glossing. The use of this package is NOT explained in the Quick Reference Guide, but the documentation is on CTAN for those that are interested. MJKD finds it handy for glossing. (https://ctan.org/pkg/leipzig?lang=en)
\usepackage{leipzig}

%Tables
% \usepackage{caption} %For table captions 

%For OT-style tableaux
\usepackage[notipa]{ot-tableau}

%highlights text with \hl{text}
\usepackage{color, soul} %note that soul sometimes eats Unicode text witouth telling you. 

%Drawing Syntax Trees
\usepackage[linguistics]{forest}

%This specifies some formatting for the forest trees to make them nicer to look at. Unfortunately this was borrowed by Michael Diercks from someone a while ago and he can't remember who. Apologies and if someone lets him know he will update this.
\forestset{
  nice nodes/.style={
    for tree={
      inner sep=0pt,
      fit=band,
    },
  },
  default preamble=nice nodes,
}

%A couple of packages that seemed to prefer being called toward the end of the preamble

%This package provides macros for typesetting SPE-style phonological rules. I've redefined the \oneof command here, so that there the rule includes a right brace. 
\usepackage{phonrule}
\renewcommand{\oneof}[2][c]{\ensuremath{
    ~\left\{
    \begin{tabular}{#1#1}#2\end{tabular}
    \right\}
    }}

%For using Greek letters outside of math mode.
% \usepackage{textgreek}


%Random, lets us use the XeLaTeX logo. Not important to the template at all.
% \usepackage{metalogo}



%%%%%%%%%%%%%%%%%%%%%%%%%%%%%%%%%%%%%%%%%%%%%
%%%%%%%%%% From Smith paper %%%%%%%%%%%%%%%%%
%%%%%%%%%%%%%%%%%%%%%%%%%%%%%%%%%%%%%%%%%%%%%

\usetikzlibrary{positioning}
\usetikzlibrary{trees}
% \usepackage{pstricks} %pstricks is banned. Use tikz instead. 
% \usepackage{pst-node}
%\usepackage{tikz}

%%%%%%%%%%%%%%%%%%%%%%%%%%%%%%%%%%%%%%%%%%%%%
%%%%%%%%%% From Gluckman paper %%%%%%%%%%%%%%
%%%%%%%%%%%%%%%%%%%%%%%%%%%%%%%%%%%%%%%%%%%%%

\usepackage{amsmath}



%%%%%%%%%%%%%%%%%%%%%%%%%%%%%%%%%%%%%%%%%%%%%
%%%%%%%%%% From Kandybowicz paper %%%%%%%%%%%
%%%%%%%%%%%%%%%%%%%%%%%%%%%%%%%%%%%%%%%%%%%%%

%These are in the .tex, but there's nothing in the "Figures" folder so I'm not sure that it's actually doing anything. 
 
% \graphicspath{ {./Figures/} }

%%%%%%%%%%%%%%%%%%%%%%%%%%%%%%%%%%%%%%%%%%%%%
%%%%%%%%%% From Matte paper %%%%%%%%%%%%%%%%%
%%%%%%%%%%%%%%%%%%%%%%%%%%%%%%%%%%%%%%%%%%%%%

%%%%%%%%%%%%%%%%%%%%%%%%%%%%%%%%%%%%%%%%%%%%%
%%%%%%%%%% From Grollemund paper %%%%%%%%%%%%%%
%%%%%%%%%%%%%%%%%%%%%%%%%%%%%%%%%%%%%%%%%%%%%

% \usepackage{lscape}

%%%%%%%%%%%%%%%%%%%%%%%%%%%%%%%%%%%%%%%%%%%%%
%%%%%%%%%% From Burns paper %%%%%%%%%%%%%%
%%%%%%%%%%%%%%%%%%%%%%%%%%%%%%%%%%%%%%%%%%%%%


% \usepackage{url}
% \urlstyle{same}

\usepackage{listings}
\lstset{basicstyle=\ttfamily,tabsize=2,breaklines=true}

%MJKD commented out - these are standard for chapters in edited volumes, I don't think we need them here in the preamble. 

% \usepackage{tabularx,multicol}
% \usepackage{langsci-basic}
% \usepackage{langsci-gb4e}
% \bibliography{localbibliography}
% \newcommand{\orcid}[1]{}
% \pagenumbering{roman}

% \IfFileExists{../localcommands.tex}{%hack to check whether this is being compiled as part of a collection or standalone
%    %%%%%%%%%%%%%%%%%%%%%%%%%%%%%%%%%%%%%%%%%%%%%
%%%%%%%%%% From LSP %%%%%%%%%%%%%%%%%%%%%%%%%
%%%%%%%%%%%%%%%%%%%%%%%%%%%%%%%%%%%%%%%%%%%%%

\usepackage{tabularx,multicol}
\usepackage{url}
\urlstyle{same}

\usepackage{langsci-optional}
\usepackage{langsci-lgr}
\usepackage{langsci-gb4e}

% NOTE THAT THE FOLLOWING PACKAGES ARE BANNED: 
% - pstricks
% - tipa
%%%%%%%%%%%%%%%%%%%%%%%%%%%%%%%%%%%%%%%%%%%%%
%%%%%%%%%% From ACAL LaTeX template %%%%%%%%%
%%%%%%%%%%%%%%%%%%%%%%%%%%%%%%%%%%%%%%%%%%%%%

%For stacking text, used here in autosegmental diagrams
\usepackage{stackengine}

%To combine rows in tables
\usepackage{multirow}

%Gives some extra formatting options, e.g. underlining/strikeout
\usepackage[normalem]{ulem}

%Use package/command below to create a double-spaced document, if you want one. Uncomment BOTH the package and the command (\doublespacing) to create a doublespaced document, or leave them as is to have a single-spaced document.
%\usepackage{setspace}
%\doublespacing 

 

%use for special OT tableaux symbols like bomb and sad face. must be loaded early on because it doesn't play well with some other packages
\usepackage{fourier-orns}

%Basic math symbols 
\usepackage{pifont}
\usepackage{amssymb}

%%%Gives shortcuts for glossing. The use of this package is NOT explained in the Quick Reference Guide, but the documentation is on CTAN for those that are interested. MJKD finds it handy for glossing. (https://ctan.org/pkg/leipzig?lang=en)
\usepackage{leipzig}

%Tables
% \usepackage{caption} %For table captions 

%For OT-style tableaux
\usepackage[notipa]{ot-tableau}

%highlights text with \hl{text}
\usepackage{color, soul} %note that soul sometimes eats Unicode text witouth telling you. 

%Drawing Syntax Trees
\usepackage[linguistics]{forest}

%This specifies some formatting for the forest trees to make them nicer to look at. Unfortunately this was borrowed by Michael Diercks from someone a while ago and he can't remember who. Apologies and if someone lets him know he will update this.
\forestset{
  nice nodes/.style={
    for tree={
      inner sep=0pt,
      fit=band,
    },
  },
  default preamble=nice nodes,
}

%A couple of packages that seemed to prefer being called toward the end of the preamble

%This package provides macros for typesetting SPE-style phonological rules. I've redefined the \oneof command here, so that there the rule includes a right brace. 
\usepackage{phonrule}
\renewcommand{\oneof}[2][c]{\ensuremath{
    ~\left\{
    \begin{tabular}{#1#1}#2\end{tabular}
    \right\}
    }}

%For using Greek letters outside of math mode.
% \usepackage{textgreek}


%Random, lets us use the XeLaTeX logo. Not important to the template at all.
% \usepackage{metalogo}



%%%%%%%%%%%%%%%%%%%%%%%%%%%%%%%%%%%%%%%%%%%%%
%%%%%%%%%% From Smith paper %%%%%%%%%%%%%%%%%
%%%%%%%%%%%%%%%%%%%%%%%%%%%%%%%%%%%%%%%%%%%%%

\usetikzlibrary{positioning}
\usetikzlibrary{trees}
% \usepackage{pstricks} %pstricks is banned. Use tikz instead. 
% \usepackage{pst-node}
%\usepackage{tikz}

%%%%%%%%%%%%%%%%%%%%%%%%%%%%%%%%%%%%%%%%%%%%%
%%%%%%%%%% From Gluckman paper %%%%%%%%%%%%%%
%%%%%%%%%%%%%%%%%%%%%%%%%%%%%%%%%%%%%%%%%%%%%

\usepackage{amsmath}



%%%%%%%%%%%%%%%%%%%%%%%%%%%%%%%%%%%%%%%%%%%%%
%%%%%%%%%% From Kandybowicz paper %%%%%%%%%%%
%%%%%%%%%%%%%%%%%%%%%%%%%%%%%%%%%%%%%%%%%%%%%

%These are in the .tex, but there's nothing in the "Figures" folder so I'm not sure that it's actually doing anything. 
 
% \graphicspath{ {./Figures/} }

%%%%%%%%%%%%%%%%%%%%%%%%%%%%%%%%%%%%%%%%%%%%%
%%%%%%%%%% From Matte paper %%%%%%%%%%%%%%%%%
%%%%%%%%%%%%%%%%%%%%%%%%%%%%%%%%%%%%%%%%%%%%%

%%%%%%%%%%%%%%%%%%%%%%%%%%%%%%%%%%%%%%%%%%%%%
%%%%%%%%%% From Grollemund paper %%%%%%%%%%%%%%
%%%%%%%%%%%%%%%%%%%%%%%%%%%%%%%%%%%%%%%%%%%%%

% \usepackage{lscape}

%%%%%%%%%%%%%%%%%%%%%%%%%%%%%%%%%%%%%%%%%%%%%
%%%%%%%%%% From Burns paper %%%%%%%%%%%%%%
%%%%%%%%%%%%%%%%%%%%%%%%%%%%%%%%%%%%%%%%%%%%%


% \usepackage{url}
% \urlstyle{same}

\usepackage{listings}
\lstset{basicstyle=\ttfamily,tabsize=2,breaklines=true}

%MJKD commented out - these are standard for chapters in edited volumes, I don't think we need them here in the preamble. 

% \usepackage{tabularx,multicol}
% \usepackage{langsci-basic}
% \usepackage{langsci-gb4e}
% \bibliography{localbibliography}
% \newcommand{\orcid}[1]{}
% \pagenumbering{roman}

% \IfFileExists{../localcommands.tex}{%hack to check whether this is being compiled as part of a collection or standalone
%    %%%%%%%%%%%%%%%%%%%%%%%%%%%%%%%%%%%%%%%%%%%%%
%%%%%%%%%% From LSP %%%%%%%%%%%%%%%%%%%%%%%%%
%%%%%%%%%%%%%%%%%%%%%%%%%%%%%%%%%%%%%%%%%%%%%

\usepackage{tabularx,multicol}
\usepackage{url}
\urlstyle{same}

\usepackage{langsci-optional}
\usepackage{langsci-lgr}
\usepackage{langsci-gb4e}

% NOTE THAT THE FOLLOWING PACKAGES ARE BANNED: 
% - pstricks
% - tipa
%%%%%%%%%%%%%%%%%%%%%%%%%%%%%%%%%%%%%%%%%%%%%
%%%%%%%%%% From ACAL LaTeX template %%%%%%%%%
%%%%%%%%%%%%%%%%%%%%%%%%%%%%%%%%%%%%%%%%%%%%%

%For stacking text, used here in autosegmental diagrams
\usepackage{stackengine}

%To combine rows in tables
\usepackage{multirow}

%Gives some extra formatting options, e.g. underlining/strikeout
\usepackage[normalem]{ulem}

%Use package/command below to create a double-spaced document, if you want one. Uncomment BOTH the package and the command (\doublespacing) to create a doublespaced document, or leave them as is to have a single-spaced document.
%\usepackage{setspace}
%\doublespacing 

 

%use for special OT tableaux symbols like bomb and sad face. must be loaded early on because it doesn't play well with some other packages
\usepackage{fourier-orns}

%Basic math symbols 
\usepackage{pifont}
\usepackage{amssymb}

%%%Gives shortcuts for glossing. The use of this package is NOT explained in the Quick Reference Guide, but the documentation is on CTAN for those that are interested. MJKD finds it handy for glossing. (https://ctan.org/pkg/leipzig?lang=en)
\usepackage{leipzig}

%Tables
% \usepackage{caption} %For table captions 

%For OT-style tableaux
\usepackage[notipa]{ot-tableau}

%highlights text with \hl{text}
\usepackage{color, soul} %note that soul sometimes eats Unicode text witouth telling you. 

%Drawing Syntax Trees
\usepackage[linguistics]{forest}

%This specifies some formatting for the forest trees to make them nicer to look at. Unfortunately this was borrowed by Michael Diercks from someone a while ago and he can't remember who. Apologies and if someone lets him know he will update this.
\forestset{
  nice nodes/.style={
    for tree={
      inner sep=0pt,
      fit=band,
    },
  },
  default preamble=nice nodes,
}

%A couple of packages that seemed to prefer being called toward the end of the preamble

%This package provides macros for typesetting SPE-style phonological rules. I've redefined the \oneof command here, so that there the rule includes a right brace. 
\usepackage{phonrule}
\renewcommand{\oneof}[2][c]{\ensuremath{
    ~\left\{
    \begin{tabular}{#1#1}#2\end{tabular}
    \right\}
    }}

%For using Greek letters outside of math mode.
% \usepackage{textgreek}


%Random, lets us use the XeLaTeX logo. Not important to the template at all.
% \usepackage{metalogo}



%%%%%%%%%%%%%%%%%%%%%%%%%%%%%%%%%%%%%%%%%%%%%
%%%%%%%%%% From Smith paper %%%%%%%%%%%%%%%%%
%%%%%%%%%%%%%%%%%%%%%%%%%%%%%%%%%%%%%%%%%%%%%

\usetikzlibrary{positioning}
\usetikzlibrary{trees}
% \usepackage{pstricks} %pstricks is banned. Use tikz instead. 
% \usepackage{pst-node}
%\usepackage{tikz}

%%%%%%%%%%%%%%%%%%%%%%%%%%%%%%%%%%%%%%%%%%%%%
%%%%%%%%%% From Gluckman paper %%%%%%%%%%%%%%
%%%%%%%%%%%%%%%%%%%%%%%%%%%%%%%%%%%%%%%%%%%%%

\usepackage{amsmath}



%%%%%%%%%%%%%%%%%%%%%%%%%%%%%%%%%%%%%%%%%%%%%
%%%%%%%%%% From Kandybowicz paper %%%%%%%%%%%
%%%%%%%%%%%%%%%%%%%%%%%%%%%%%%%%%%%%%%%%%%%%%

%These are in the .tex, but there's nothing in the "Figures" folder so I'm not sure that it's actually doing anything. 
 
% \graphicspath{ {./Figures/} }

%%%%%%%%%%%%%%%%%%%%%%%%%%%%%%%%%%%%%%%%%%%%%
%%%%%%%%%% From Matte paper %%%%%%%%%%%%%%%%%
%%%%%%%%%%%%%%%%%%%%%%%%%%%%%%%%%%%%%%%%%%%%%

%%%%%%%%%%%%%%%%%%%%%%%%%%%%%%%%%%%%%%%%%%%%%
%%%%%%%%%% From Grollemund paper %%%%%%%%%%%%%%
%%%%%%%%%%%%%%%%%%%%%%%%%%%%%%%%%%%%%%%%%%%%%

% \usepackage{lscape}

%%%%%%%%%%%%%%%%%%%%%%%%%%%%%%%%%%%%%%%%%%%%%
%%%%%%%%%% From Burns paper %%%%%%%%%%%%%%
%%%%%%%%%%%%%%%%%%%%%%%%%%%%%%%%%%%%%%%%%%%%%


% \usepackage{url}
% \urlstyle{same}

\usepackage{listings}
\lstset{basicstyle=\ttfamily,tabsize=2,breaklines=true}

%MJKD commented out - these are standard for chapters in edited volumes, I don't think we need them here in the preamble. 

% \usepackage{tabularx,multicol}
% \usepackage{langsci-basic}
% \usepackage{langsci-gb4e}
% \bibliography{localbibliography}
% \newcommand{\orcid}[1]{}
% \pagenumbering{roman}

% \IfFileExists{../localcommands.tex}{%hack to check whether this is being compiled as part of a collection or standalone
%    \input{../localpackages}
%    \input{../localcommands}
%    \input{../localhyphenation}
%     \bibliography{localbibliography}
%     \togglepaper[23]
% }{}

% %custom footer for preprints
% \papernote{\scriptsize\normalfont
%     Change author.
%     Change title.
%     To appear in:
%     Change Volume Editor.  
%     Change volume title.  
%     Berlin: Language Science Press. [preliminary page numbering]
% }



%%%%%%%%%%%%%%%%%%%%%%%%%%%%%%%%%%%%%%%%%%%%%
%%%%%%%%%% From Feibig paper %%%%%%%%%%%%%%
%%%%%%%%%%%%%%%%%%%%%%%%%%%%%%%%%%%%%%%%%%%%%

\usepackage{tikz-qtree}
\usepackage{tikz-qtree-compat}

%%%%%%%%%%%%%%%%%%%%%%%%%%%%%%%%%%%%%%%%%%%%%
%%%%%%%%%% From Olson paper %%%%%%%%%%%%%%
%%%%%%%%%%%%%%%%%%%%%%%%%%%%%%%%%%%%%%%%%%%%%

%MJKD - TIPA is forbidden lol. Commenting out for now but we'll have to address this in the Olson paper

%\usepackage[tone]{tipa}

%%%%%%%%%%%%%%%%%%%%%%%%%%%%%%%%%%%%%%%%%%%%%
%%%%%%%%%% From Caragine paper %%%%%%%%%%%%%%
%%%%%%%%%%%%%%%%%%%%%%%%%%%%%%%%%%%%%%%%%%%%%

\usepackage{placeins}
% \usepackage{graphicx}
% \usepackage{float} %Overleaf suggested this as a solution


%%%%%%%%%%%%%%%%%%%%%%%%%%%%%%%%%%%%%%%%%%%%%
%%%%%%%%%% From Kieffer paper %%%%%%%%%%%%%%
%%%%%%%%%%%%%%%%%%%%%%%%%%%%%%%%%%%%%%%%%%%%%

\usepackage{array}

%%%%%%%%%%%%%%%%%%%%%%%%%%%%%%%%%%%%%%%%%%%%%
%%%%%%%%%% From Payne paper %%%%%%%%%%%%%%
%%%%%%%%%%%%%%%%%%%%%%%%%%%%%%%%%%%%%%%%%%%%%

\usepackage{enumerate}
\usepackage{array,ragged2e}
\usepackage{pgfplots}


%%%%%%%%%%%%%%%%%%%%%%%%%%%%%%%%%%%%%%%%%%%%%
%%%%%%%%%% From Diercks paper %%%%%%%%%%%%%%
%%%%%%%%%%%%%%%%%%%%%%%%%%%%%%%%%%%%%%%%%%%%%

% \usepackage{cgloss}



%%%%%%%%%%%%%%%%%%%%%%%%%%%%%%%%%%%%%%%%%%%%%
%%%%%%%%%% From Li paper %%%%%%%%%%%%%%
%%%%%%%%%%%%%%%%%%%%%%%%%%%%%%%%%%%%%%%%%%%%%



%%%%%%%%%%%%%%%%%%%%%%%%%%%%%%%%%%%%%%%%%%%%%
%%%%%%%%%% From Myers paper %%%%%%%%%%%%%%
%%%%%%%%%%%%%%%%%%%%%%%%%%%%%%%%%%%%%%%%%%%%%



\usepackage{tikz-qtree}
\usepackage{tikz-qtree-compat}


%Package that makes glossing slightly easier.
\RequirePackage{leipzig}
%\RequirePackage{mathtools} % for \mathrlap

% % % Shortcuts (various borrowed from Jason Zentz)
\RequirePackage{xspace}
\xspaceaddexceptions{]\}}
\usepackage{langsci-textipa}
\usepackage{langsci-branding}
\usepackage{tabto}

%    %%%%%%%%%%%%%%%%%%%%%%%%%%%%%%%%%%%%%%%%%%%%%
%%%%%%%%%% From LSP %%%%%%%%%%%%%%%%%%%%%%%%%
%%%%%%%%%%%%%%%%%%%%%%%%%%%%%%%%%%%%%%%%%%%%%


% \newcommand*{\orcid}{}


\makeatletter
\let\thetitle\@title
\let\theauthor\@author
\makeatother

\newcommand{\togglepaper}[1][0]{
   \bibliography{../localbibliography}
   \papernote{\scriptsize\normalfont
     \theauthor.
     \titleTemp.
     To appear in:
     E. Di Tor \& Herr Rausgeberin (ed.).
     Booktitle in localcommands.tex.
     Berlin: Language Science Press. [preliminary page numbering]
   }
   \pagenumbering{roman}
   \setcounter{chapter}{#1}
   \addtocounter{chapter}{-1}
}

\newbool{bookcompile}
\booltrue{bookcompile}
\newcommand{\bookorchapter}[2]{\ifbool{bookcompile}{#1}{#2}}


%%%%%%%%%%%%%%%%%%%%%%%%%%%%%%%%%%%%%%%%%%%%%
%%%%%%%%%% From ACAL LaTeX template %%%%%%%%%
%%%%%%%%%%%%%%%%%%%%%%%%%%%%%%%%%%%%%%%%%%%%%


\usepackage{tikz}

\newcommand{\circled}[1]{\begin{tikzpicture}[baseline=(word.base)]
\node[draw, rounded corners, text height=8pt, text depth=2pt, inner sep=2pt, outer sep=0pt, use as bounding box] (word) {#1};
\end{tikzpicture}
}


%%%%%%%%%%%%%%%%%%%%%%%%%%%%%%%%%%%%%%%%%%%%%%
%%%%%%%%%%%%%%%%%%%%%%%%%%%%%%%%%%%%%%%%%%%%%%

% Useful Ling Shortcuts


%This makes the \emptyset command be a nicer one
\let\oldemptyset\emptyset
\let\emptyset\varnothing
\newcommand{\nothing}{$\emptyset$}

%Not all of these are explained in the Quick Reference Guide, but they are here bc they are relevant to some of our students. 
\newcommand{\xbar}{\ensuremath{'}\xspace}
\newcommand{\xhead}{\ensuremath{^\circ}\xspace}
\newcommand{\Lb}[1]{$\text{[}_{\text{#1}}$ } %A more convenient left bracket
\newcommand{\Rb}[1]{$\text{]}_{\text{#1}}$ } %A more convenient left bracket
\newcommand{\gap}{\underline{\hspace{1.2em}}}
\newcommand{\vP}{\emph{v}P}
\newcommand{\lilv}{\emph{v}}
\newcommand{\Abar}{A$'$-} %A more convenient A-bar notation
\newcommand{\ph}{$\varphi$\xspace} %A more convenient phi
\newcommand{\pro}{\emph{pro}\xspace}
\newcommand{\subs}[1]{\textsubscript{#1}} %A more convenient subscript
\newcommand{\spells}{$\Longleftrightarrow$} %spellout arrow for morph spellout rules
\newcommand{\tr}[1]{\textit{t}\textsubscript{\textit{#1}}} %easy traces with subscript
\newcommand{\supers}[1]{\textsuperscript{#1}}

%These are borrowed/modified from Andrew Mackenzie's ling-macros package. We have copied them in here (instead of just using the package itself) to avoid a minor clash that was generating an error.


% Abbreviations for glossing, based on Leipzig
\newleipzig{hab}{hab}{habitual}
\newleipzig{rem}{rem}{remote}
\newleipzig{sm}{sm}{subject marker}
\newleipzig{t}{t}{tense}
\newleipzig{aa}{aa}{anti-agreement}
\newleipzig{pron}{pron}{pronoun}
\newleipzig{rec}{rec}{recent}
\newleipzig{om}{om}{object marker}
%\newleipzig{ipfv}{ipfv}{imperfective}
\newleipzig{asp}{asp}{aspect}
\newleipzig{lk}{lk}{linker}
\newleipzig{pcl}{pcl}{particle}
\newleipzig{stat}{stat}{stative}
\newleipzig{ints}{ints}{intensive}
\newleipzig{ascl}{ascl}{assertive subject clitic}
\newleipzig{nascl}{nascl}{non-assertive subject clitic}
\newleipzig{ta}{ta}{tense and/or aspect}
\newleipzig{assoc}{assoc}{associative marker}
\newleipzig{hon}{hon}{honorific}
%\newleipzig{whprt}{wh}{\wh particle}
\newleipzig{sa}{sa}{subject agreement}
\newleipzig{conj}{conj}{conjunction}
%\newleipzig{loc}{loc}{locative}
\newleipzig{expl}{expl}{expletive}
\newleipzig{rcm}{rcm}{reciprocal marker}
\newleipzig{pers}{pers}{persistive}
\newleipzig{fv}{fv}{final vowel}
%\newleipzig{}{}{} %this is just to copy for when I want to add more


%%%%%%%%%%%%%%%%%%%%%%%%%%%%%%%%%%%%%%%%%%%%%%%%%%%%%
%%%%%%%%%%%%%%%%%%%%%%%%%%%%%%%%%%%%%%%%%%%%%%%%%%%%%


%%%%%%%%%%%%%%%%%%%%%%%%%%%%%%%%%%%%%%%%%%%%%
%%%%%%%%%% From Gluckman paper %%%%%%%%%%%%%%
%%%%%%%%%%%%%%%%%%%%%%%%%%%%%%%%%%%%%%%%%%%%%

\newcommand{\pcite}[1]{\citeauthor{#1}'s (\citeyear{#1})}


%%%%%%%%%%%%%%%%%%%%%%%%%%%%%%%%%%%%%%%%%%%%%
%%%%%%%%%% From Myers paper %%%%%%%%%%%%%%
%%%%%%%%%%%%%%%%%%%%%%%%%%%%%%%%%%%%%%%%%%%%%

%% hspace over width of something without showing it
\newcommand*{\hspaceThis}[1]{\settowidth{\LSPTmp}{#1}\hspace*{\LSPTmp}}
% use \hphantom{} for this


%%%%%%%%%%%%%%%%%%%%%%%%%%%%%%%%%%%%%%%%%%%%%
%%%%%%%%%% From Matte paper %%%%%%%%%%%%%%%%%
%%%%%%%%%%%%%%%%%%%%%%%%%%%%%%%%%%%%%%%%%%%%%

%mjkd commented this out in case it's breaking something right now. 
% \newcounter{xlistex}
% \newcommand{\footnoteex}{\stepcounter{xlistex}(\roman{xlistex})}

\newleipzig{sp}{sp}{speaker} %a new leipzig glossing command


%%%%%%%%%%%%%%%%%%%%%%%%%%%%%%%%%%%%%%%%%%%%%
%%%%%%%%%% From GamFranich paper %%%%%%%%%%%%%%
%%%%%%%%%%%%%%%%%%%%%%%%%%%%%%%%%%%%%%%%%%%%%

%MJKD commented these out - both will be provided by the overall LSP book template 

% \bibliography{localbibliography}
% \newcommand{\orcid}[1]{}

%%%%%%%%%%%%%%%%%%%%%%%%%%%%%%%%%%%%%%%%%%%%%
%%%%%%%%%% From Diercks paper %%%%%%%%%%%%%%
%%%%%%%%%%%%%%%%%%%%%%%%%%%%%%%%%%%%%%%%%%%%%

\newcommand{\abar}{$\bar{\text{A}}$\xspace} % this one is different to the \Abar command in Mike's shortcuts doc
\newcommand{\topFeat}{[\textsc{top}]\xspace}


%%%%%%%%%%%%%%%%%%%%%%%%%%%%%%%%%%%%%%%%%%%%%
%%%%%%%%%% From Kinchsular paper %%%%%%%%%%%%%%
%%%%%%%%%%%%%%%%%%%%%%%%%%%%%%%%%%%%%%%%%%%%%

%%%%%%%%%%%%%%%%%%%%%%%%%%%%%%%%%%%%%%%%%%%%%
%%%%%%%%%% From Chen paper %%%%%%%%%%%%%%
%%%%%%%%%%%%%%%%%%%%%%%%%%%%%%%%%%%%%%%%%%%%%

\newcounter{savedex}
\makeatletter
\newcommand*{\saveex}{\setcounter{savedex}{\value{\@xnumctr}}}
\newcommand*{\resumeex}{\setcounter{\@xnumctr}{\value{savedex}}}
\makeatother


%    \input{../localhyphenation}
%     \bibliography{localbibliography}
%     \togglepaper[23]
% }{}

% %custom footer for preprints
% \papernote{\scriptsize\normalfont
%     Change author.
%     Change title.
%     To appear in:
%     Change Volume Editor.  
%     Change volume title.  
%     Berlin: Language Science Press. [preliminary page numbering]
% }



%%%%%%%%%%%%%%%%%%%%%%%%%%%%%%%%%%%%%%%%%%%%%
%%%%%%%%%% From Feibig paper %%%%%%%%%%%%%%
%%%%%%%%%%%%%%%%%%%%%%%%%%%%%%%%%%%%%%%%%%%%%

\usepackage{tikz-qtree}
\usepackage{tikz-qtree-compat}

%%%%%%%%%%%%%%%%%%%%%%%%%%%%%%%%%%%%%%%%%%%%%
%%%%%%%%%% From Olson paper %%%%%%%%%%%%%%
%%%%%%%%%%%%%%%%%%%%%%%%%%%%%%%%%%%%%%%%%%%%%

%MJKD - TIPA is forbidden lol. Commenting out for now but we'll have to address this in the Olson paper

%\usepackage[tone]{tipa}

%%%%%%%%%%%%%%%%%%%%%%%%%%%%%%%%%%%%%%%%%%%%%
%%%%%%%%%% From Caragine paper %%%%%%%%%%%%%%
%%%%%%%%%%%%%%%%%%%%%%%%%%%%%%%%%%%%%%%%%%%%%

\usepackage{placeins}
% \usepackage{graphicx}
% \usepackage{float} %Overleaf suggested this as a solution


%%%%%%%%%%%%%%%%%%%%%%%%%%%%%%%%%%%%%%%%%%%%%
%%%%%%%%%% From Kieffer paper %%%%%%%%%%%%%%
%%%%%%%%%%%%%%%%%%%%%%%%%%%%%%%%%%%%%%%%%%%%%

\usepackage{array}

%%%%%%%%%%%%%%%%%%%%%%%%%%%%%%%%%%%%%%%%%%%%%
%%%%%%%%%% From Payne paper %%%%%%%%%%%%%%
%%%%%%%%%%%%%%%%%%%%%%%%%%%%%%%%%%%%%%%%%%%%%

\usepackage{enumerate}
\usepackage{array,ragged2e}
\usepackage{pgfplots}


%%%%%%%%%%%%%%%%%%%%%%%%%%%%%%%%%%%%%%%%%%%%%
%%%%%%%%%% From Diercks paper %%%%%%%%%%%%%%
%%%%%%%%%%%%%%%%%%%%%%%%%%%%%%%%%%%%%%%%%%%%%

% \usepackage{cgloss}



%%%%%%%%%%%%%%%%%%%%%%%%%%%%%%%%%%%%%%%%%%%%%
%%%%%%%%%% From Li paper %%%%%%%%%%%%%%
%%%%%%%%%%%%%%%%%%%%%%%%%%%%%%%%%%%%%%%%%%%%%



%%%%%%%%%%%%%%%%%%%%%%%%%%%%%%%%%%%%%%%%%%%%%
%%%%%%%%%% From Myers paper %%%%%%%%%%%%%%
%%%%%%%%%%%%%%%%%%%%%%%%%%%%%%%%%%%%%%%%%%%%%



\usepackage{tikz-qtree}
\usepackage{tikz-qtree-compat}


%Package that makes glossing slightly easier.
\RequirePackage{leipzig}
%\RequirePackage{mathtools} % for \mathrlap

% % % Shortcuts (various borrowed from Jason Zentz)
\RequirePackage{xspace}
\xspaceaddexceptions{]\}}
\usepackage{langsci-textipa}
\usepackage{langsci-branding}
\usepackage{tabto}

%    %%%%%%%%%%%%%%%%%%%%%%%%%%%%%%%%%%%%%%%%%%%%%
%%%%%%%%%% From LSP %%%%%%%%%%%%%%%%%%%%%%%%%
%%%%%%%%%%%%%%%%%%%%%%%%%%%%%%%%%%%%%%%%%%%%%


% \newcommand*{\orcid}{}


\makeatletter
\let\thetitle\@title
\let\theauthor\@author
\makeatother

\newcommand{\togglepaper}[1][0]{
   \bibliography{../localbibliography}
   \papernote{\scriptsize\normalfont
     \theauthor.
     \titleTemp.
     To appear in:
     E. Di Tor \& Herr Rausgeberin (ed.).
     Booktitle in localcommands.tex.
     Berlin: Language Science Press. [preliminary page numbering]
   }
   \pagenumbering{roman}
   \setcounter{chapter}{#1}
   \addtocounter{chapter}{-1}
}

\newbool{bookcompile}
\booltrue{bookcompile}
\newcommand{\bookorchapter}[2]{\ifbool{bookcompile}{#1}{#2}}


%%%%%%%%%%%%%%%%%%%%%%%%%%%%%%%%%%%%%%%%%%%%%
%%%%%%%%%% From ACAL LaTeX template %%%%%%%%%
%%%%%%%%%%%%%%%%%%%%%%%%%%%%%%%%%%%%%%%%%%%%%


\usepackage{tikz}

\newcommand{\circled}[1]{\begin{tikzpicture}[baseline=(word.base)]
\node[draw, rounded corners, text height=8pt, text depth=2pt, inner sep=2pt, outer sep=0pt, use as bounding box] (word) {#1};
\end{tikzpicture}
}


%%%%%%%%%%%%%%%%%%%%%%%%%%%%%%%%%%%%%%%%%%%%%%
%%%%%%%%%%%%%%%%%%%%%%%%%%%%%%%%%%%%%%%%%%%%%%

% Useful Ling Shortcuts


%This makes the \emptyset command be a nicer one
\let\oldemptyset\emptyset
\let\emptyset\varnothing
\newcommand{\nothing}{$\emptyset$}

%Not all of these are explained in the Quick Reference Guide, but they are here bc they are relevant to some of our students. 
\newcommand{\xbar}{\ensuremath{'}\xspace}
\newcommand{\xhead}{\ensuremath{^\circ}\xspace}
\newcommand{\Lb}[1]{$\text{[}_{\text{#1}}$ } %A more convenient left bracket
\newcommand{\Rb}[1]{$\text{]}_{\text{#1}}$ } %A more convenient left bracket
\newcommand{\gap}{\underline{\hspace{1.2em}}}
\newcommand{\vP}{\emph{v}P}
\newcommand{\lilv}{\emph{v}}
\newcommand{\Abar}{A$'$-} %A more convenient A-bar notation
\newcommand{\ph}{$\varphi$\xspace} %A more convenient phi
\newcommand{\pro}{\emph{pro}\xspace}
\newcommand{\subs}[1]{\textsubscript{#1}} %A more convenient subscript
\newcommand{\spells}{$\Longleftrightarrow$} %spellout arrow for morph spellout rules
\newcommand{\tr}[1]{\textit{t}\textsubscript{\textit{#1}}} %easy traces with subscript
\newcommand{\supers}[1]{\textsuperscript{#1}}

%These are borrowed/modified from Andrew Mackenzie's ling-macros package. We have copied them in here (instead of just using the package itself) to avoid a minor clash that was generating an error.


% Abbreviations for glossing, based on Leipzig
\newleipzig{hab}{hab}{habitual}
\newleipzig{rem}{rem}{remote}
\newleipzig{sm}{sm}{subject marker}
\newleipzig{t}{t}{tense}
\newleipzig{aa}{aa}{anti-agreement}
\newleipzig{pron}{pron}{pronoun}
\newleipzig{rec}{rec}{recent}
\newleipzig{om}{om}{object marker}
%\newleipzig{ipfv}{ipfv}{imperfective}
\newleipzig{asp}{asp}{aspect}
\newleipzig{lk}{lk}{linker}
\newleipzig{pcl}{pcl}{particle}
\newleipzig{stat}{stat}{stative}
\newleipzig{ints}{ints}{intensive}
\newleipzig{ascl}{ascl}{assertive subject clitic}
\newleipzig{nascl}{nascl}{non-assertive subject clitic}
\newleipzig{ta}{ta}{tense and/or aspect}
\newleipzig{assoc}{assoc}{associative marker}
\newleipzig{hon}{hon}{honorific}
%\newleipzig{whprt}{wh}{\wh particle}
\newleipzig{sa}{sa}{subject agreement}
\newleipzig{conj}{conj}{conjunction}
%\newleipzig{loc}{loc}{locative}
\newleipzig{expl}{expl}{expletive}
\newleipzig{rcm}{rcm}{reciprocal marker}
\newleipzig{pers}{pers}{persistive}
\newleipzig{fv}{fv}{final vowel}
%\newleipzig{}{}{} %this is just to copy for when I want to add more


%%%%%%%%%%%%%%%%%%%%%%%%%%%%%%%%%%%%%%%%%%%%%%%%%%%%%
%%%%%%%%%%%%%%%%%%%%%%%%%%%%%%%%%%%%%%%%%%%%%%%%%%%%%


%%%%%%%%%%%%%%%%%%%%%%%%%%%%%%%%%%%%%%%%%%%%%
%%%%%%%%%% From Gluckman paper %%%%%%%%%%%%%%
%%%%%%%%%%%%%%%%%%%%%%%%%%%%%%%%%%%%%%%%%%%%%

\newcommand{\pcite}[1]{\citeauthor{#1}'s (\citeyear{#1})}


%%%%%%%%%%%%%%%%%%%%%%%%%%%%%%%%%%%%%%%%%%%%%
%%%%%%%%%% From Myers paper %%%%%%%%%%%%%%
%%%%%%%%%%%%%%%%%%%%%%%%%%%%%%%%%%%%%%%%%%%%%

%% hspace over width of something without showing it
\newcommand*{\hspaceThis}[1]{\settowidth{\LSPTmp}{#1}\hspace*{\LSPTmp}}
% use \hphantom{} for this


%%%%%%%%%%%%%%%%%%%%%%%%%%%%%%%%%%%%%%%%%%%%%
%%%%%%%%%% From Matte paper %%%%%%%%%%%%%%%%%
%%%%%%%%%%%%%%%%%%%%%%%%%%%%%%%%%%%%%%%%%%%%%

%mjkd commented this out in case it's breaking something right now. 
% \newcounter{xlistex}
% \newcommand{\footnoteex}{\stepcounter{xlistex}(\roman{xlistex})}

\newleipzig{sp}{sp}{speaker} %a new leipzig glossing command


%%%%%%%%%%%%%%%%%%%%%%%%%%%%%%%%%%%%%%%%%%%%%
%%%%%%%%%% From GamFranich paper %%%%%%%%%%%%%%
%%%%%%%%%%%%%%%%%%%%%%%%%%%%%%%%%%%%%%%%%%%%%

%MJKD commented these out - both will be provided by the overall LSP book template 

% \bibliography{localbibliography}
% \newcommand{\orcid}[1]{}

%%%%%%%%%%%%%%%%%%%%%%%%%%%%%%%%%%%%%%%%%%%%%
%%%%%%%%%% From Diercks paper %%%%%%%%%%%%%%
%%%%%%%%%%%%%%%%%%%%%%%%%%%%%%%%%%%%%%%%%%%%%

\newcommand{\abar}{$\bar{\text{A}}$\xspace} % this one is different to the \Abar command in Mike's shortcuts doc
\newcommand{\topFeat}{[\textsc{top}]\xspace}


%%%%%%%%%%%%%%%%%%%%%%%%%%%%%%%%%%%%%%%%%%%%%
%%%%%%%%%% From Kinchsular paper %%%%%%%%%%%%%%
%%%%%%%%%%%%%%%%%%%%%%%%%%%%%%%%%%%%%%%%%%%%%

%%%%%%%%%%%%%%%%%%%%%%%%%%%%%%%%%%%%%%%%%%%%%
%%%%%%%%%% From Chen paper %%%%%%%%%%%%%%
%%%%%%%%%%%%%%%%%%%%%%%%%%%%%%%%%%%%%%%%%%%%%

\newcounter{savedex}
\makeatletter
\newcommand*{\saveex}{\setcounter{savedex}{\value{\@xnumctr}}}
\newcommand*{\resumeex}{\setcounter{\@xnumctr}{\value{savedex}}}
\makeatother


%    \input{../localhyphenation}
%     \bibliography{localbibliography}
%     \togglepaper[23]
% }{}

% %custom footer for preprints
% \papernote{\scriptsize\normalfont
%     Change author.
%     Change title.
%     To appear in:
%     Change Volume Editor.  
%     Change volume title.  
%     Berlin: Language Science Press. [preliminary page numbering]
% }



%%%%%%%%%%%%%%%%%%%%%%%%%%%%%%%%%%%%%%%%%%%%%
%%%%%%%%%% From Feibig paper %%%%%%%%%%%%%%
%%%%%%%%%%%%%%%%%%%%%%%%%%%%%%%%%%%%%%%%%%%%%

\usepackage{tikz-qtree}
\usepackage{tikz-qtree-compat}

%%%%%%%%%%%%%%%%%%%%%%%%%%%%%%%%%%%%%%%%%%%%%
%%%%%%%%%% From Olson paper %%%%%%%%%%%%%%
%%%%%%%%%%%%%%%%%%%%%%%%%%%%%%%%%%%%%%%%%%%%%

%MJKD - TIPA is forbidden lol. Commenting out for now but we'll have to address this in the Olson paper

%\usepackage[tone]{tipa}

%%%%%%%%%%%%%%%%%%%%%%%%%%%%%%%%%%%%%%%%%%%%%
%%%%%%%%%% From Caragine paper %%%%%%%%%%%%%%
%%%%%%%%%%%%%%%%%%%%%%%%%%%%%%%%%%%%%%%%%%%%%

\usepackage{placeins}
% \usepackage{graphicx}
% \usepackage{float} %Overleaf suggested this as a solution


%%%%%%%%%%%%%%%%%%%%%%%%%%%%%%%%%%%%%%%%%%%%%
%%%%%%%%%% From Kieffer paper %%%%%%%%%%%%%%
%%%%%%%%%%%%%%%%%%%%%%%%%%%%%%%%%%%%%%%%%%%%%

\usepackage{array}

%%%%%%%%%%%%%%%%%%%%%%%%%%%%%%%%%%%%%%%%%%%%%
%%%%%%%%%% From Payne paper %%%%%%%%%%%%%%
%%%%%%%%%%%%%%%%%%%%%%%%%%%%%%%%%%%%%%%%%%%%%

\usepackage{enumerate}
\usepackage{array,ragged2e}
\usepackage{pgfplots}


%%%%%%%%%%%%%%%%%%%%%%%%%%%%%%%%%%%%%%%%%%%%%
%%%%%%%%%% From Diercks paper %%%%%%%%%%%%%%
%%%%%%%%%%%%%%%%%%%%%%%%%%%%%%%%%%%%%%%%%%%%%

% \usepackage{cgloss}



%%%%%%%%%%%%%%%%%%%%%%%%%%%%%%%%%%%%%%%%%%%%%
%%%%%%%%%% From Li paper %%%%%%%%%%%%%%
%%%%%%%%%%%%%%%%%%%%%%%%%%%%%%%%%%%%%%%%%%%%%



%%%%%%%%%%%%%%%%%%%%%%%%%%%%%%%%%%%%%%%%%%%%%
%%%%%%%%%% From Myers paper %%%%%%%%%%%%%%
%%%%%%%%%%%%%%%%%%%%%%%%%%%%%%%%%%%%%%%%%%%%%



\usepackage{tikz-qtree}
\usepackage{tikz-qtree-compat}


%Package that makes glossing slightly easier.
\RequirePackage{leipzig}
%\RequirePackage{mathtools} % for \mathrlap

% % % Shortcuts (various borrowed from Jason Zentz)
\RequirePackage{xspace}
\xspaceaddexceptions{]\}}
\usepackage{langsci-textipa}
\usepackage{langsci-branding}
\usepackage{tabto}

   %%%%%%%%%%%%%%%%%%%%%%%%%%%%%%%%%%%%%%%%%%%%%
%%%%%%%%%% From LSP %%%%%%%%%%%%%%%%%%%%%%%%%
%%%%%%%%%%%%%%%%%%%%%%%%%%%%%%%%%%%%%%%%%%%%%


% \newcommand*{\orcid}{}


\makeatletter
\let\thetitle\@title
\let\theauthor\@author
\makeatother

\newcommand{\togglepaper}[1][0]{
   \bibliography{../localbibliography}
   \papernote{\scriptsize\normalfont
     \theauthor.
     \titleTemp.
     To appear in:
     E. Di Tor \& Herr Rausgeberin (ed.).
     Booktitle in localcommands.tex.
     Berlin: Language Science Press. [preliminary page numbering]
   }
   \pagenumbering{roman}
   \setcounter{chapter}{#1}
   \addtocounter{chapter}{-1}
}

\newbool{bookcompile}
\booltrue{bookcompile}
\newcommand{\bookorchapter}[2]{\ifbool{bookcompile}{#1}{#2}}


%%%%%%%%%%%%%%%%%%%%%%%%%%%%%%%%%%%%%%%%%%%%%
%%%%%%%%%% From ACAL LaTeX template %%%%%%%%%
%%%%%%%%%%%%%%%%%%%%%%%%%%%%%%%%%%%%%%%%%%%%%


\usepackage{tikz}

\newcommand{\circled}[1]{\begin{tikzpicture}[baseline=(word.base)]
\node[draw, rounded corners, text height=8pt, text depth=2pt, inner sep=2pt, outer sep=0pt, use as bounding box] (word) {#1};
\end{tikzpicture}
}


%%%%%%%%%%%%%%%%%%%%%%%%%%%%%%%%%%%%%%%%%%%%%%
%%%%%%%%%%%%%%%%%%%%%%%%%%%%%%%%%%%%%%%%%%%%%%

% Useful Ling Shortcuts


%This makes the \emptyset command be a nicer one
\let\oldemptyset\emptyset
\let\emptyset\varnothing
\newcommand{\nothing}{$\emptyset$}

%Not all of these are explained in the Quick Reference Guide, but they are here bc they are relevant to some of our students. 
\newcommand{\xbar}{\ensuremath{'}\xspace}
\newcommand{\xhead}{\ensuremath{^\circ}\xspace}
\newcommand{\Lb}[1]{$\text{[}_{\text{#1}}$ } %A more convenient left bracket
\newcommand{\Rb}[1]{$\text{]}_{\text{#1}}$ } %A more convenient left bracket
\newcommand{\gap}{\underline{\hspace{1.2em}}}
\newcommand{\vP}{\emph{v}P}
\newcommand{\lilv}{\emph{v}}
\newcommand{\Abar}{A$'$-} %A more convenient A-bar notation
\newcommand{\ph}{$\varphi$\xspace} %A more convenient phi
\newcommand{\pro}{\emph{pro}\xspace}
\newcommand{\subs}[1]{\textsubscript{#1}} %A more convenient subscript
\newcommand{\spells}{$\Longleftrightarrow$} %spellout arrow for morph spellout rules
\newcommand{\tr}[1]{\textit{t}\textsubscript{\textit{#1}}} %easy traces with subscript
\newcommand{\supers}[1]{\textsuperscript{#1}}

%These are borrowed/modified from Andrew Mackenzie's ling-macros package. We have copied them in here (instead of just using the package itself) to avoid a minor clash that was generating an error.


% Abbreviations for glossing, based on Leipzig
\newleipzig{hab}{hab}{habitual}
\newleipzig{rem}{rem}{remote}
\newleipzig{sm}{sm}{subject marker}
\newleipzig{t}{t}{tense}
\newleipzig{aa}{aa}{anti-agreement}
\newleipzig{pron}{pron}{pronoun}
\newleipzig{rec}{rec}{recent}
\newleipzig{om}{om}{object marker}
%\newleipzig{ipfv}{ipfv}{imperfective}
\newleipzig{asp}{asp}{aspect}
\newleipzig{lk}{lk}{linker}
\newleipzig{pcl}{pcl}{particle}
\newleipzig{stat}{stat}{stative}
\newleipzig{ints}{ints}{intensive}
\newleipzig{ascl}{ascl}{assertive subject clitic}
\newleipzig{nascl}{nascl}{non-assertive subject clitic}
\newleipzig{ta}{ta}{tense and/or aspect}
\newleipzig{assoc}{assoc}{associative marker}
\newleipzig{hon}{hon}{honorific}
%\newleipzig{whprt}{wh}{\wh particle}
\newleipzig{sa}{sa}{subject agreement}
\newleipzig{conj}{conj}{conjunction}
%\newleipzig{loc}{loc}{locative}
\newleipzig{expl}{expl}{expletive}
\newleipzig{rcm}{rcm}{reciprocal marker}
\newleipzig{pers}{pers}{persistive}
\newleipzig{fv}{fv}{final vowel}
%\newleipzig{}{}{} %this is just to copy for when I want to add more


%%%%%%%%%%%%%%%%%%%%%%%%%%%%%%%%%%%%%%%%%%%%%%%%%%%%%
%%%%%%%%%%%%%%%%%%%%%%%%%%%%%%%%%%%%%%%%%%%%%%%%%%%%%


%%%%%%%%%%%%%%%%%%%%%%%%%%%%%%%%%%%%%%%%%%%%%
%%%%%%%%%% From Gluckman paper %%%%%%%%%%%%%%
%%%%%%%%%%%%%%%%%%%%%%%%%%%%%%%%%%%%%%%%%%%%%

\newcommand{\pcite}[1]{\citeauthor{#1}'s (\citeyear{#1})}


%%%%%%%%%%%%%%%%%%%%%%%%%%%%%%%%%%%%%%%%%%%%%
%%%%%%%%%% From Myers paper %%%%%%%%%%%%%%
%%%%%%%%%%%%%%%%%%%%%%%%%%%%%%%%%%%%%%%%%%%%%

%% hspace over width of something without showing it
\newcommand*{\hspaceThis}[1]{\settowidth{\LSPTmp}{#1}\hspace*{\LSPTmp}}
% use \hphantom{} for this


%%%%%%%%%%%%%%%%%%%%%%%%%%%%%%%%%%%%%%%%%%%%%
%%%%%%%%%% From Matte paper %%%%%%%%%%%%%%%%%
%%%%%%%%%%%%%%%%%%%%%%%%%%%%%%%%%%%%%%%%%%%%%

%mjkd commented this out in case it's breaking something right now. 
% \newcounter{xlistex}
% \newcommand{\footnoteex}{\stepcounter{xlistex}(\roman{xlistex})}

\newleipzig{sp}{sp}{speaker} %a new leipzig glossing command


%%%%%%%%%%%%%%%%%%%%%%%%%%%%%%%%%%%%%%%%%%%%%
%%%%%%%%%% From GamFranich paper %%%%%%%%%%%%%%
%%%%%%%%%%%%%%%%%%%%%%%%%%%%%%%%%%%%%%%%%%%%%

%MJKD commented these out - both will be provided by the overall LSP book template 

% \bibliography{localbibliography}
% \newcommand{\orcid}[1]{}

%%%%%%%%%%%%%%%%%%%%%%%%%%%%%%%%%%%%%%%%%%%%%
%%%%%%%%%% From Diercks paper %%%%%%%%%%%%%%
%%%%%%%%%%%%%%%%%%%%%%%%%%%%%%%%%%%%%%%%%%%%%

\newcommand{\abar}{$\bar{\text{A}}$\xspace} % this one is different to the \Abar command in Mike's shortcuts doc
\newcommand{\topFeat}{[\textsc{top}]\xspace}


%%%%%%%%%%%%%%%%%%%%%%%%%%%%%%%%%%%%%%%%%%%%%
%%%%%%%%%% From Kinchsular paper %%%%%%%%%%%%%%
%%%%%%%%%%%%%%%%%%%%%%%%%%%%%%%%%%%%%%%%%%%%%

%%%%%%%%%%%%%%%%%%%%%%%%%%%%%%%%%%%%%%%%%%%%%
%%%%%%%%%% From Chen paper %%%%%%%%%%%%%%
%%%%%%%%%%%%%%%%%%%%%%%%%%%%%%%%%%%%%%%%%%%%%

\newcounter{savedex}
\makeatletter
\newcommand*{\saveex}{\setcounter{savedex}{\value{\@xnumctr}}}
\newcommand*{\resumeex}{\setcounter{\@xnumctr}{\value{savedex}}}
\makeatother


   \input{../localhyphenation}
   \boolfalse{bookcompile}
   \togglepaper[23]%%chapternumber
}{}

\begin{document}
\maketitle
\section{Introduction} \label{sect:grollemundbrownrule:introduction}

Sawabantu Bantu languages are located along the Cameroonian, Equatorial Guinean and Gabonese coasts (\figref{fig:grollemundbrownrule:languagearea}) and generally include Guthrie’s (\citealt{guthriecomparativebantu}, later updated by \citealt{maho2009}) A10, A20 and A30 languages. Cameroon is well-known for its significant cultural and linguistic diversity with approximately 250 spoken languages \citep{eberhardetal2019} belonging to three different language families: Niger-Congo (190 languages including the Bantoid and Bantu languages), Nilo-Saharan (2 languages) and Afro-Asiatic (55 languages). Moreover, it is in southern Cameroon that the Bantu Expansion started 5,000 years ago with Bantu-speaking-populations migrating southward through Cameroon, Equatorial Guinea and Gabon expanding later across Central, Eastern and Southern Africa \citep{bostoen2007, grollemundetal2015, koileetal2022}.

Several studies have been conducted on Sawabantu languages highlighting their relatedness, with A10-20-30 languages (minus Manenguba A15 and Bubi A31) forming a valid genetic group. \citet{ebobisse1989} proposed first a dialectometrical study of A20 and A30 Sawabantu languages for which he calculated the percentages of linguistic proximity based on the study of lexicon. His results divided the Sawabantu clade into two subgroups with A20 languages on one hand and A30 languages on the other hand. In 2015, \citet{ebobisse2015} furthered his study of Sawabantu languages by providing a detailed synchronic study of the same languages with a focus on phonology, morphology and lexicon confirming the two subgroups identified in 1989. The addition of A10 languages (minus Manenguba A15) sometimes called “Oroko” \citep{friesen2002} to the Sawabantu clade was made by \citet[170]{nurseandphilippson2003} who established a classification of 80 Bantu languages based on the study of non-lexical data (30 phonological and morphological features). The study of lexicon but also morphological criteria such as the shape of the nominal stem that is generally CVCV in A10-20-30 as opposed to Manenguba A15 \citep[6-7]{philippson2022} were decisive in their classification. Based on these findings, \citet{philippson2022} proposed to borrow the geographical term “Sawabantu” from \citet{ebobisse2015} to refer to these A10-20-30 languages that have a genetic value.

The purpose of this paper is to examine Sawabantu languages and reconstruct their putative ancestor Proto-Sawabantu (PSB) by the application of the Comparative Method (CM). PSB will be analyzed and compared with other proto-languages reconstructed in this region such as Proto-Manenguba (PM) reconstructed by \citet{hedinger1987} that includes A15 languages spoken near the Bantu cradle. The comparison between PSB and PM reconstructions reveals two different evolutionary trajectories with PSB being more similar to Proto-Bantu (PB) in certain aspects as opposed to PM confirming therefore the differences observed between Manenguba A15 and Sawabantu A10-A20-A30 languages. PSB reconstructions will also be compared to PB reconstructions taken from BLR3 \citep{bastinetal2002} in order to identify any diachronic phonological innovations that will shed new light on the peculiar evolution of the languages spoken in this region. 

\section{Data and methods} \label{sect:grollemundbrownrule:dataandmethods}

\subsection{Languages} \label{sect:grollemundbrownrule:languages} 

Based on the study of lexicon \citep{ebobisse1989, ebobisse2015} and non-lexical data such as phonological or morphological features \citep[70]{nurseandphilippson2003}, the Sawabantu clade generally comprises between 20--27 Bantu languages belonging to Guthrie’s zones A10, A20 and A30 (e.g., \tabref{tab:grollemundbrownrule:languages}) with the exception of the A15 languages called Manenguba and A31 Bubi spoken on Bioko Island. 

In order to start our comparative work, we selected the same A20 and A30 Sawabantu languages provided in \citet{ebobisse2015} study. \citet{ebobisse2015} divided the Sawabantu languages into three groups based on geographical locations: South (Banoho, Batanga, Yasa, Bapuku), Central (Duala, Bidiman, Oli, Pongo, Mongo, Mulimba) and Southwest (Mboke, Mokpe, Isu, Kole, Bubia) but his classification based on the study of lexicon published in 1989 showed that Central and Southwest were more closely related. We will use in this paper the same geographical distinctions established in \citet{ebobisse1989} and \citet{ebobisse2015} with the mention of South, Central, Southwest with the addition of the North region (“Oroko”) represented in our study by the Kundu language (A10) that we have extracted from \citet{friesen2002} in order to get a more representative sample of the Sawabantu languages.

In his monumental study, \citet{ebobisse2015} collected, compared and analyzed 610 words for 15 languages. For the purpose of this study, we created an Excel database containing all 610 words for the same 15 languages with the addition of Kundu, an A10 language. Words from Ebobissé were translated from German to English. In addition to collecting the individual words, we searched for corresponding PB words drawn from BLR 3 \citep{bastinetal2002} for comparison purposes.

\begin{table}
\caption{Sawabantu languages in \citet{ebobisse2015} with selected languages for this study in italics with Maho’s codes (\citeyear{maho2009}) and ISO codes when available (some languages do not have ISO code because they are either understudied/unknown or considered as some variants of another language).}
\label{tab:grollemundbrownrule:languages}
 \begin{tabular}{llll}
  \lsptoprule
          Region  & Language & Maho (2009) & ISO 639-3\\
  \midrule
  South  &   \textit{Banoho (BAN)} & A32a &   \\
    &   \textit{Batanga (BAT)} & A32 &  bnm  \\
    &  \textit{Yasa (YAS)} & A33a &  yko  \\
    &   \textit{Bapuku (BAP)} & A32b &  bng  \\
    &   Kombe & A33b &  nui  \\
    &   Benga & A34 &  bng  \\
  Central  &   \textit{Duala (DUA)} & A24 & dua  \\
    &   \textit{Bidiman (BID)} & A241 &   \\
    &   \textit{Oli (OLI)} & A25 &   \\
    &  \textit{ Pongo (PON)} & A26 &   \\
    &   \textit{Mongo (MON)} & A261 &   \\
    &   \textit{Mulimba (MUL)} & A27 & mzd  \\
Southwest  &   \textit{Mboko (MBO)} & A21 & mdu  \\
    &   \textit{Mokpe (KPE)} & A22 & bri  \\
    &  \textit{Isu (ISU)} & A23 & srz  \\
    &   \textit{Kole (KOL)} & A231 & kme \\
    &   \textit{Bubia (BUB)} & A221 & bbx  \\
North  &   Londo  & A11 &  \\
    &   Ngolo  & A111 &   \\
    &   Bima  & A112 &   \\
    &   Batanga  & A113 &   \\
    &   Bakoko  & A114 &   \\
    &   Londo  & A115 &   \\
    &   Lue  & A12 &   \\
    &   Mbonge  & A121 &   \\
    &   \textit{Kundu}  & A122 &   \\
    &   Ekombe  & A123 &   \\          
  \lspbottomrule
 \end{tabular}
\end{table}

\begin{figure}
    \centering
%     \includegraphics[width=0.5\linewidth]{figures/Grollemundfigure01map.png}
    \includegraphics[width=0.8\linewidth]{figures/grollemund.pdf}
    \caption{Approximate location of Sawabantu languages (in the red rectangle)}
   \label{fig:grollemundbrownrule:languagearea}
\end{figure}

\subsection{Lexicon-based classifications of Sawabantu languages} \label{sect:browngrollemundrule:previousstudies}

Bantu languages spoken in southern Cameroon belong to the North-Western (NW) genetic group \citep{bastinetal1999, grollemund2012, currieetal2013, grollemundetal2015, koileetal2022} that usually includes all the A languages (Guthrie’s (\citeyear{guthriecomparativebantu}) referential zones) generally with B10-20-30 languages spoken in Gabon. \citet{bastinetal1999}, for example, include in their NW group A languages (without Mbam-Bubi composed of A31-some A40-50-60 languages) and B10-20-30 languages. According to Bastin et al. 1999’s classification, Mbam-Bubi languages are the first group to diverge followed by the NW. \citet{grollemundetal2015} shows that NW was composed of several clades that generally includes: Mbam-Bubi (+Jarawan), A10-20-30-40+A70-80-90 (labeled “NW Cameroon”), some B20+A80-90, remaining B20 and B10-30 (labeled “NW Gabon”). The latest phylogenetic classification established by \citet{koileetal2022} proposes a different result wherein NW is composed of a sole node that includes only A languages and some of B20; whereas the remaining B20 and B10-30 languages are classified into a different clade. To sum up, Sawabantu languages belong to the first group of Bantu languages to diverge after the Mbam-Bubi group. 

A closer look at the Bantu trees with a focus on the NW clades (\figref{fig:grollemundbrownrule:comparison}) established by Grollemund et al. (2015) and Koile et al. (2022) confirms that Mbam languages (A40-50-60) are the first Bantu languages to diverge followed by a group that includes A10-20-30-40-70 (in \citealt{grollemundetal2015}) or only A10-20-30 (in \citealt{koileetal2022}). In both trees, Sawabantu languages split after the Mbam languages and form a clade with Manenguba (A15) languages confirming their relatedness. In addition, both classifications highlight a noticeable division amongst Sawabantu languages with a first group composed of A24-A27-A33a-A34 (corresponding to Ebobissé’s South and Central groups) and a second group composed of A22-A122 languages (Southwest) linked with two A10 languages Londo and Mbonge (North). This observation contradicts Ebobissé’s (\citeyear{ebobisse1989}) results where Central and Southwest are more closely related but confirms Nurse and Philippson’s (\citeyear{nurseandphilippson2003}) classification with the inclusion of North Sawabantu languages within the Sawabantu group. Finally, we observe that this last subgroup (A22-A122-A11-A121) is linked to the Manenguba languages (A15) as a coordinated branch. 

\begin{figure}
    \centering
    \includegraphics[width=0.8\linewidth]{figures/Grollemundfigure02trees.png}
    \caption{Comparison of the two NW groups in Grollemund et al. (2015) and Koile et al. (2022)}
   \label{fig:grollemundbrownrule:comparison}
\end{figure}

\subsection{Method} \label{sect:browngrollemundrule:methods}

In order to reconstruct PSB, we applied the Comparative Method (CM) and its series of steps. This upstream approach utilizes the systematic comparison of cognate words in related languages to reconstruct the linguistic properties of the putative ancestor with a great degree of certainty \citep{campbell1998}. In order to successfully apply the CM, several steps must be completed. First, the identification of cognates - words that are similar in form and meaning that are assumed to have been inherited from a common ancestral source - is paramount. We studied each meaning to identify cognate sets\footnote{Because of the space restriction, we are not able to present our data with our 600 comparative series. However, many examples are presented in this paper with approximately 100 different comparative series for the languages studied.}. Most of the words studied yielded between one to three cognate sets (e.g., \tabref{tab:grollemundbrownrule:cognates} `two', `house' and `surpass') whereas, in rare cases, meanings would yield as many as five different cognate sets (e.g., \tabref{tab:grollemundbrownrule:cognates} `yesterday'). 

\begin{table}
\caption{Examples of cognacy judgement, numbers refer to cognate sets. The capital letters refer to the geographical denominations with “S” for South, “C” for Central, “SW” for Southwest and “N” for North. Empty cells mark missing data.}
\label{tab:grollemundbrownrule:cognates}
 \begin{tabular}{llllll}
  \lsptoprule
          {{Meanings}} & {{two}} & {{house}} & {{breast}} & {{surpass}} & {{yesterday}} \\
  \midrule
  BAN - S                    & ɓé-ɓà 1             & ndáɓò 1                 & ɗi-ɓɛ́ɛ 1                & i-ɓála 1                 & vákɛ́ 1                    \\
BAT - S                    & ɓé-ɓà 1             & ndáɓò 1                 & ɗi-ɓɛ́ɛ 1                & i-ɓála 1                 & váhɛ́ 1                    \\
YAS - S                    & ɓí-ɓà 1             &                         & i-ɓɛ́nɛ 1                & e-ɓálɛ 1                 &                             \\
BAP - S                    & ɓé-ɓà 1             & ndáɓò 1                 & i-ɓɛ́ɛ 1                 & i-ɓála 1                 &                             \\
DUA - C                    & ɓé-ɓàá 1           & ndáɓò 1                 & di-ɓɛ́ɛ 1                & búka 2                   &                             \\
BID - C                    & ɓí-ɓàá 1           & ndáɓò 1                 & ɓɛ́ɛ 1                   & búa 2                    & ɓwéná 2                   \\
OLI - C                    & ɓí-ɓàá 1           & ndáɓò 1                 & ɓɛ́ɛ 1                   & búa 2                    & ɓwéá 2                    \\
PON - C                    & ɓé-ɓàá 1           & ndáɓò 1                 & i-ɓɛ́ɛ 1                 & i-búka 2                 & mbánā 2                   \\
MON - C                    & ɓé-ɓàá 1           & ndáɓò 1                 & di-ɓɛ́ɛ 1                & i-búka 2                 & kɛ́kɛ 3                     \\
MUL - C                    & ɓé-ɓàá 1           & ndáɓò 1                 & di-ɓɛ́ɛ 1                &                           & cɛ́cɛ 3                     \\
MBO - SW                   & βé-βàá 1           & ndáwò 1                 & li-βɛ́ɛ 1                & laáka 3                  &                             \\
KPE - SW                   & βé-βàá 1           & ndáwò 1                 & li-βɛ́ɛ 1                & laáka 3                  & ŋmɛ́ɛlɛ́ɛ 4                 \\
ISU - SW                   & βé-ɓàá 1           & ndáwò 1                 & li-βɛ́ɛ 1                & yaáka 3                  & kɛ́ɛkɛ́ɛ 5                  \\
KOL - SW                   & bé-bàá 1           & ndábò 1                 & di-bɛ́ɛ 1                & yaáka 3                  & kɛ́kɛ́ɛ 5                   \\
BUB - SW                   & βé-βàá 1           & ndáwò 1                 & li-βɛ́ɛ 1                & láka 3                   & ŋméelɛ́ɛ 4                 \\
Kundu - N                  & βe-βɛ 1               & ndáβo 1                 & di-bɛ̂ 1                 &                           & jana 2 \\                    
   \lspbottomrule
 \end{tabular}
\end{table}

Once all the cognate sets were identified, we proceeded to the next step which consisted of the alignment of all the cognate Bantu roots in C1, V1, C2 and V2 positions in order to list all the sound correspondences present in the data. The sound correspondences (study of the phonological similarities between the languages studied) were then organized to facilitate the identification of regular sound correspondences. If a sound correspondence occurs several times, it suggests that the sound correspondence is regular which implies that the sounds could have either been inherited from proto-sounds located at an ancestral level or innovated. For example, if in all the Sawabantu languages we found the sound /m/ in C1 and C2 positions, it is reasonable to assume that the PSB had *m in the same positions. The establishment of regular sound correspondences is then used to reconstruct ancestral sounds. Finally, once the regular sound changes have been understood via the study of regular sound correspondences, historical linguists can proceed with the reconstruction of words. 

\section{Results}
\largerpage
\subsection{Comparative tables} \label{sect:browngrollemundrule:comparativetables}

In this section, we present the consonant and vowel sound correspondences as well as the reconstructed sounds labeled “PSB”. The reconstructions made are primarily based on noun and verb stems which generally present a CVCV structure for nouns and CVC structure for verbs in Sawabantu languages.

\subsubsection{Vowels in C1 and C2 positions} \label{sect:browngrollemundrule:vowels}

\begin{table}[b]
\caption{Reconstruction of PSB Vowels}
\label{tab:grollemundbrownrule:psbvowels}
 \begin{tabular}{lrrrrrrr}
  \lsptoprule
 \textbf{PSB} & \textbf{*i} & \textbf{*e} & \textbf{*ɛ} & \textbf{*a} & \textbf{*ɔ} & \textbf{*o} & \textbf{*u} \\
   \midrule
BAN - S      & i           & e           & ɛ           & a           & ɔ           & o           & u           \\
BAT - S      & i           & e           & ɛ           & a           & ɔ           & o           & u           \\
YAS - S      & i           & e           & ɛ           & a           & ɔ           & o           & u           \\
BAP - S      & i           & e           & ɛ           & a           & ɔ           & o           & u           \\
DUA - C      & i           & e           & ɛ           & a           & ɔ           & o           & u           \\
BID - C      & i           & e           & ɛ           & a           & ɔ           & o           & u           \\
OLI - C      & i           & e           & ɛ           & a           & ɔ           & o           & u           \\
PON - C      & i           & e           & ɛ           & a           & ɔ           & o           & u           \\
MON - C      & i           & e           & ɛ           & a           & ɔ           & o           & u           \\
MUL - C      & i           & e           & ɛ           & a           & ɔ           & o           & u           \\
MBO - SW     & i           & e           & ɛ           & a           & ɔ           & o           & u           \\
KPE - SW     & i           & e           & ɛ           & a           & ɔ           & o           & u           \\
ISU - SW     & i           & e           & ɛ           & a           & ɔ           & o           & u           \\
KOL - SW     & i           & e           & ɛ           & a           & ɔ           & o           & u           \\
BUB - SW     & i           & e           & ɛ           & a           & ɔ           & o           & u           \\
Kundu - N    & i           & e           & ɛ           & a           & ɔ           & o           & u          \\
   \lspbottomrule
 \end{tabular}
\end{table}


We reconstructed a system with seven vowels at the PSB level. Examining regular sound correspondences for vowels (presented in \tabref{tab:grollemundbrownrule:psbvowels}) shows that vowels are highly regular in Sawabantu languages. This regularity contrasts with Manenguba languages where vowels have undergone a number of changes since PM such as assimilation, raising or lowering which complicate the reconstruction according to \citet[44-45]{hedinger1987}.  \citet[45-46]{hedinger1987} also reconstructs long vowels for PM that are generally the result of two distinct processes: (i) loss of the PB *C2 in Manenguba languages has created CVV roots (PB *bèdú `colanut' > PM *bèé) and (ii) compensatory lengthening in CVC roots after C2 was deleted (PM *fiɲ `pus' > /-híí/ in Myenge). 


\begin{table}[b]
\caption{Examples for the reflexes of PSB vowels in V1 and V2 positions. Empty cells represent either missing data or non-cognate words.}
\label{tab:grollemundbrownrule:vowelexamples}
\fittable{
 \begin{tabular}{llllllll}
  \lsptoprule
\textbf{PSB}              & \textbf{*i$_{12}$}           & \textbf{*e$_{12}$}          & \textbf{*ɛ$_{12}$}           & \textbf{*a$_{12}$}            & \textbf{*ɔ$_{12}$}          & \textbf{*o$_{12}$}          & \textbf{*u$_{12}$}            \\
{{meaning}} & {{cloud}} & {{tree}} & {{stand}} & {{divide}} & {{knee}} & {{skin}} & {{turtle}} \\
\midrule
BAN - S                   & ɗi-vîndî                & yelé                  & i-tɛ́mɛ                 & i-kaɓa                   & ɗì-ɓɔ́ŋgɔ́            & ŋ̀-hoɓo                & kúɗù                   \\
BAT - S                   & ɗi-vîndî                & yelé                  & i-tɛ́mɛ                 & i-kaɓa                   & ɗì-ɓɔ́ŋgɔ́            & ŋ̀-hoɓo                & kúɗù                   \\
YAS - S                   & i-vîndó                 & e-yilé                & e-tɛ́mɛ                 & e-kaɓa                   & ì-ɓɔ́ŋgɔ́             & mo-koɓo                & kúɗù                   \\
BAP - S                   & ɗi-vîndî                & yelé                  & i-tɛ́mɛ                 & i-kaɓa                   & ì-ɓɔ́ŋgɔ́             & ŋ̀-koɓo                & kúɗù                   \\
DUA - C                   & di-wîndî                & ɓw-elé                & tɛ́mɛ                   & j-aɓa                    & dì-ɓɔ́ŋgɔ́            & e-yoɓo                 & wúdù                   \\
BID - C                   & findîfindî              & ɓw-elé                & tɛ́mɛ                   & y-aɓa                    & ɗ-òŋgò               & o-ɓo                   & wúlù                   \\
OLI - C                   & vîndîvîndî              & ɓw-elé                & tɛ́mɛ                   & y-aɓa                    & ɓɔ́ŋ                   & o-ɓo                   & wúlù                   \\
PON - C                   & i-wîndîwîndî            & ɓw-elé                & i-tɛ́mɛ                 & ij-aɓa                   & ì-ɓɔ́ŋgɔ́             & e-yoɓo                 & wúdù                   \\
MON - C                   & di-wîndîwîndî           & ɓw-elé                & i-tɛ́mɛ                 & -kaɓanɛ                  & dì-ɓɔ́ŋgɔ́            & e-koɓo                 & kúdù                   \\
MUL - C                   & di-îndî                 & ɓw-elé                & i-tɛ́mɛ                 & i-kaɓa                   & dì-ɓɔ́ŋgɔ́            & e-yoɓo                 & kúdù                   \\
MBO - SW                  & li-wîndó                & eelé                  & li-tɛ́mɛ                & li-βaá                  & lì-wɔ́ŋgɔ́ŋgɔ́        & e-kowo                 & è-kúlɛ̀                \\
KPE - SW                  & li-wîndî                & ɡb-eé                 & li-tɛ́mɛ                & li-βaá                  & lì-wɔ́ŋgɔ́ŋgɔ́        & yowo                   & kúlù                   \\
ISU - SW                  & li-wîndîwîndî           & ɡb-elé                & i-tɛ́mɛ                 & i-kaβa                   & lì-wɔ́ŋgɔ́ŋgɔ́        & e-kowo                 & è-kúlɛ̀kùlɛ̀          \\
KOL - SW                  & di-wîndî                & ɡb-eré                & i-tɛ́mɛ                 &                          & dì-bɔ́ŋgɔ́            &                        & è-kúrɛ̀                \\
BUB - SW                  & li-îndó                 & eelé                  & li-tɛ́mɛ                & li-kaβa                  & lì-wɔ́ŋgɔ́ŋgɔ́        & m-owo                  & è-kúlɛ̀kùlɛ̀          \\
Kundu - N                 &                         & bo-le                  &                         & -kaβa                    & di-bɔŋgɔ               & e-kobo                 & ku                           \\ 
   \lspbottomrule
 \end{tabular}
 }
\end{table}


The situation is less complicated in Sawabantu languages where we mostly have CVCV roots for noun stems. There are some languages where non-contrastive long vowels are attested. Generally, they are also the results of the loss of C2 (such as PSB *k) leading to the emergence of CVV roots. In such cases, long vowels or VV sequences can emerge such as in PSB *toko `spoon' > /too/ or PSB *sako `army ant' > /sao/ in some Central Sawabantu languages (where PSB *k2 > Ø). Therefore, because this phenomenon is caused by the deletion of a C2 that is still attested in the remaining Sawabantu languages (as opposed to PM where the PB *C2 has completely disappeared in all Manenguba daughter languages), we chose to not reconstruct long vowels.

The system reconstructed for PSB is a typical PB 7-vowel system with *ɪ and *ʊ for Bantu languages spoken in Central Africa \citep{guthriecomparativebantu, meeussen1967, bastinetal2002}. The sole disagreement that existed with the reconstruction of PB vowels was regarding the height of the high vowels *i, *ɪ, *u, *ʊ \citep{bastinetal2002} vs *į, *i, *ų, *u \citep{guthriecomparativebantu, meeussen1967}. While \citet{guthriecomparativebantu} and \citet{meeussen1967} propose that some PB vowels were super close (marked with a cedilla), \citet{bastinetal2002} choose to have a contrast between tense and lax high PB vowels. Some modern-day Bantu languages present the tense/lax opposition ([ɪ] and [ʊ]) whereas other languages, such as Sawabantu languages, have mid vowels ([e] and [o]) instead.

\subsubsection{Consonants} \label{sect:browngrollemundrule:consonants}

We present in this section the consonants reconstructed for PSB. When needed, a distinction between C1 (initial consonant) and C2 (intervocalic consonant) is made (e.g., *k$_1$ vs *k$_2$ in \tabref{tab:grollemundbrownrule:consonants}). First, we present the voiceless plosives in \tabref{tab:grollemundbrownrule:consonants}. 

\begin{table}
\caption{Sound correspondences for unvoiced plosives. When two reflexes are presented, the first one is the most frequent. Bold characters mark sound changes from PSB to the daughter languages.}
\label{tab:grollemundbrownrule:consonants}
\fittable{
 \begin{tabular}{llllllll}
  \lsptoprule
\textbf{PSB} & \textbf{*p$_{12}$} & \textbf{*t$_{12}$} & \textbf{*k$_1$} & \textbf{*k$_1$ / \_Vk$_2$} & \textbf{*k$_1$ /\_e} & \textbf{*k$_1$ /\_wV} & \textbf{*k$_2$} \\
\midrule
BAN - S      & p             & t             & \textbf{h}$\sim$k    & k                    & k$\sim$h         & (w)                & \textbf{h}$\sim$k    \\
BAT - S      & p             & t             & \textbf{h}$\sim$k    & k                    & k$\sim$h         & (w)                & \textbf{h}$\sim$k    \\
YAS - S      & p             & t             & k            & k                    & \textbf{c}                 & (w)                & k            \\
BAP - S      & p             & t             & k            & k                    & \textbf{s}                 & (w)                & k            \\
DUA - C      & p             & t             & \textbf{Ø}$\sim$k    & k                    & Ø                 & k(w)               & \textbf{Ø}$\sim$k    \\
BID - C      & p             & t / \textbf{r} (*t$_2$)   & \textbf{Ø}$\sim$k    & k                    & Ø                 & k(w)               & \textbf{ɣ$\sim$Ø}    \\
OLI - C      & p             & t / \textbf{r} (*t$_2$)   & \textbf{Ø}$\sim$k    & k                    & Ø                 & k(w)               & \textbf{ɣ$\sim$Ø}    \\
PON - C      & p             & t             & \textbf{Ø}$\sim$k    & k                    & Ø                 & k(w)               & \textbf{Ø}$\sim$k    \\
MON - C      & p             & t             & k            & k                    & k                 & k(w)               & k            \\
MUL - C      & p             & t             & \textbf{Ø}$\sim$k    & k                    & \textbf{c}                 & k(w)               & \textbf{Ø}$\sim$k    \\
MBO - SW     & \textbf{ɸ}             & t             & k            & k                    & k                 & \textbf{kp}                 & k            \\
KPE - SW     & \textbf{ɸ}             & t             & \textbf{Ø}$\sim$k    & k                    & \textbf{Ø}$\sim$k         & \textbf{kp}                 & \textbf{Ø}$\sim$k    \\
ISU - SW     & \textbf{ɸ}             & t             & k            & k                    & k                 & \textbf{kp}                 & k            \\
KOL - SW     & \textbf{ɸ}             & t             & \textbf{Ø}$\sim$k    & k                    & \textbf{Ø}$\sim$k         & \textbf{kp}                 & \textbf{Ø}$\sim$k    \\
BUB - SW     & \textbf{ɸ}             & t             & \textbf{Ø}$\sim$k    & k                    & \textbf{Ø}$\sim$k         & \textbf{kp}                 & \textbf{Ø}$\sim$k    \\
Kundu - N    & \textbf{ɸ}             & t             & k            & k                    & k                 & \textbf{kp}                 & k            \\
   \lspbottomrule
 \end{tabular}
 }
\end{table}

For the voiceless plosives presented in \tabref{tab:grollemundbrownrule:consonants} (examples are presented in \tabref{tab:grollemundbrownrule:psbpexamples}), we reconstruct *p in C1 and C2 positions for the set p/ɸ with South and Central Sawabantu languages displaying /p/ whereas Southwest and North Sawabantu languages have the fricative /ɸ/. Based on the principles of majority wins used in CM and directionality (lenition is more common than fortition, (\citealt{bybeeandeasterday2019}), *p is the logical choice. 

\begin{table}
\caption{Examples for the reflexes of PSB *p in C1 and C2 positions. We can notice some voicing occurring in Bidiman with *p$_2$ > /ɓ/ in the examples presented here (“lip”, “touch”), this phenomenon is irregular as we can find *p$_2$ > /p/ in between vowels (“sweep”: [papa] in Bidiman).}
\label{tab:grollemundbrownrule:psbpexamples}

\fittable{
 \begin{tabular}{lllllllll}
  \lsptoprule
\textbf{PSB}              & \textbf{*p$_1$}            & \textbf{*p$_1$}            & \textbf{*p$_1$}            & \textbf{*p$_1$}           & \textbf{*p$_{12}$}         & \textbf{*p$_2$}            & \textbf{*p$_2$}             & *p$_2$                    \\
{{meaning}} & {{viper}} & {{chalk}} & {{fruit}} & {{bury}} & {{lip}} & {{touch}} & {{armpit}} & {{wind}} \\
\midrule
BAN - S                   & \textbf{p}ée                    & \textbf{p}ɛ́mbɛ́                 & e-ɓumá                  & i-\textbf{p}uɗɛ                 & è-\textbf{p}ò\textbf{p}ò             &                         & jájá\textbf{p}ì                & mpu\textbf{p}ɛ́                 \\
BAT - S                   & \textbf{p}ée                    & \textbf{p}ɛ́mbɛ́                 & e-\textbf{p}umá                  & i-\textbf{p}uɗɛ                 & è-\textbf{p}ò\textbf{p}ò             & i-tá\textbf{p}a                  & jájá\textbf{p}ì                & mpu\textbf{p}ɛ́                 \\
YAS - S                   & \textbf{p}éyi                   & \textbf{p}ɛ́mbɛ́                 & e-ɓumá                  &                        &                       &                         & ɲɔ́ɲɔ́\textbf{p}ì                & mpu\textbf{p}ɛ́                 \\
BAP - S                   & \textbf{p}ée                    & \textbf{p}ɛ́mbɛ́                 & e-ɓumá                  & i-\textbf{p}uɗɛ                 & è-\textbf{p}ò\textbf{p}ò             &                         & ɲá\textbf{ɸ}ɛ̀                   & mpu\textbf{p}ɛ́                 \\
DUA - C                   & \textbf{p}é                     & \textbf{p}ɛ́mbɛ́                 & e-\textbf{p}umá                  & \textbf{p}ulɛ                   & è-\textbf{p}ò\textbf{p}ò             & tá\textbf{p}a                    & è-yá\textbf{p}ì                  &                        \\
BID - C                   & \textbf{p}é                     & \textbf{p}ɛ́mbɛ́                 & \textbf{p}umá                    &                        & \textbf{p}ò\textbf{ɓ}ò                & tá\textbf{ɓ}a                    &                          &                        \\
OLI - C                   & \textbf{p}é                     & \textbf{p}ɛ́mbɛ́                 & \textbf{p}umá                    &                        & \textbf{p}ò\textbf{p}ò                & tá\textbf{p}a                    &                          &                        \\
PON - C                   & \textbf{p}é                     & \textbf{p}ɛ́mbɛ́                 & e-\textbf{p}umá                  & i-\textbf{p}ulɛ                 & è-\textbf{p}ò\textbf{p}ò             & i-tá\textbf{p}a                  & è-yá\textbf{p}ì                  &                        \\
MON - C                   & \textbf{p}é                     & \textbf{p}ɛ́mbɛ́                 & e-\textbf{p}umá                  & i-\textbf{p}ulɛ                 & è-\textbf{p}ò\textbf{p}ò             & i-tá\textbf{p}a                  & è-yá\textbf{p}ì                  &                        \\
MUL - C                   & \textbf{p}é                     & \textbf{p}ɛ́mbɛ́                 & e-\textbf{p}umá                  & i-\textbf{p}uɗɛ                 & è-\textbf{p}ò\textbf{p}ò             & i-tá\textbf{p}a                  & yàá\textbf{p}ì                    & mu-pu\textbf{p}é                \\
MBO - SW                  & \textbf{ɸ}ée                    & \textbf{ɸ}ɛ́mbɛ́                 & e-\textbf{ɸ}umá                  &                        &                       &                         & ɲá\textbf{ɸ}ɛ̀                   &                        \\
KPE - SW                  & \textbf{ɸ}ée                    & \textbf{ɸ}ɛ́mbɛ́                 & e-\textbf{ɸ}umá                  &                        &                       &                         & ɲáɸà\textbf{ɸ}ɛ̀                &                        \\
ISU - SW                  & \textbf{ɸ}é                     & \textbf{ɸ}ɛ́mbɛ́                 & e-\textbf{ɸ}umá                  & i-\textbf{ɸ}ulɛ                 & è-\textbf{ɸ}ò\textbf{ɸ}ò             & i-tá\textbf{ɸ}à                & ɲáɸà\textbf{ɸ}è                &                        \\
    KOL - SW                  & \textbf{ɸ}ée                    & \textbf{ɸ}ɛ́mbɛ́                 & e-\textbf{ɸ}umá                  & i-\textbf{ɸ}urɛ                 & è-\textbf{ɸ}ò\textbf{ɸ}ò             & i-tá\textbf{ɸ}à                & ɲáɸà\textbf{ɸ}ì                &                        \\
BUB - SW                  & \textbf{ɸ}ée                    & \textbf{ɸ}ɛ́mbɛ́                 & e-\textbf{ɸ}umá                  &                        &                       &                         & ɲáɸà\textbf{ɸ}è                &                        \\
Kundu - N                 &                         &                         & e-\textbf{ɸ}uma                  & \textbf{ɸ}ura                   &                       &                         &                          & di-ɸɛ\textbf{ɸ}ɛ                                  \\ 
   \lspbottomrule
 \end{tabular}
 }
\end{table}

We chose to reconstruct *t in C1 and C2 positions for the set t/r (Oli and Bidiman) because all the Sawabantu languages display /t/ in C1 and C2 positions except for Bidiman and Oli where *t$_2$ undergoes lenition (rhotacism) only in C2 position (e.g., \tabref{tab:grollemundbrownrule:psbtexamples} for examples). 


\begin{table}[p]
\caption{Examples for the reflexes of PSB *t in C1 and C2 positions. Only one example displays *t$_1$ > /r/ in Oli and Bidiman (e.g., “ear”). }
\label{tab:grollemundbrownrule:psbtexamples}
\fittable{
 \begin{tabular}{llllllll}
  \lsptoprule
\textbf{PSB}              & \textbf{*t$_1$}           & \textbf{*t$_1$}          & \textbf{*t$_1$}            & \textbf{*t$_1$}           & \textbf{*t$_{12}$}           & \textbf{*t$_2$}          & \textbf{*t$_2$}           \\
{{meaning}} & {{five}} & {{ear}} & {{drill}} & {{send}} & {{smoke}} & {{bag}} & {{hide}} \\
\midrule
BAN - S                   & ɓé-\textbf{t}áno                & ì-\textbf{t}óì              & i-\textbf{t}úɓa                  & i-\textbf{t}ila                 & í-\textbf{t}ù\textbf{t}ù               & ŋ̀-kù\textbf{t}à             & i-hu\textbf{t}a                 \\
BAT - S                   & ɓé-\textbf{t}áno                & ì-\textbf{t}óì              & i-\textbf{t}úɓa                  & i-\textbf{t}ila                 & í-\textbf{t}ù\textbf{t}ù               & ŋ̀-kù\textbf{t}à             & i-hu\textbf{t}a                 \\
YAS - S                   & ɓî-\textbf{t}áno                & ɗì-\textbf{l}ɔ́ɔ̀             & e-\textbf{t}úɓa                  & e-\textbf{t}ila                 & ví-\textbf{t}ù\textbf{t}ù              & mò-kù\textbf{t}à            &                        \\
BAP - S                   & ɓé-\textbf{t}áno                & ɗì-\textbf{t}ɔ́               & i-\textbf{t}úɓa                  & i-\textbf{t}ila                 & í-\textbf{t}ù\textbf{t}ù               &                       & i-ku\textbf{t}a                 \\
DUA - C                   & ɓé-\textbf{t}ánu                & \textbf{t}óì                 & \textbf{t}úɓa                    & \textbf{t}ila                   & í-\textbf{t}ù\textbf{t}ù               & mù-kù\textbf{t}à            & wu\textbf{t}a                   \\
BID - C                   & bî-\textbf{t}áan                & dì-\textbf{r}ɔ́               & \textbf{t}úɓa                    & \textbf{t}ila                   & mí-\textbf{t}ù\textbf{r}ù              & ŋ̀-kù\textbf{r}à             & wu\textbf{r}a                   \\
OLI - C                   & bî-\textbf{t}áan                & dì-\textbf{r}ó               & \textbf{t}úɓa                    & \textbf{t}ila                   & mí-\textbf{t}ù\textbf{r}ù              & ŋ̀-kù\textbf{r}à             & wu\textbf{r}a                   \\
PON - C                   & ɓé-\textbf{t}ánu                & ì-\textbf{t}óì              & i-\textbf{t}úɓa                  & i-\textbf{t}ila                 & í-\textbf{t}ù\textbf{t}ù               & mù-kù\textbf{t}à            & i-wu\textbf{t}a                 \\
MON - C                   & ɓé-\textbf{t}ánu                & ì-\textbf{t}óì              & i-\textbf{t}úɓa                  & i-\textbf{t}ila                 & í-\textbf{t}ù\textbf{t}ù               & mù-kù\textbf{t}à            & i-ku\textbf{t}a                 \\
MUL - C                   & ɓé-\textbf{t}áan                & ì-\textbf{t}óì              & i-\textbf{t}úɓa                  & i-\textbf{t}ila                 & í-\textbf{t}ù\textbf{t}ù               & mù-kù\textbf{t}à            & i-ku\textbf{t}a                 \\
MBO - SW                  & βé-\textbf{t}áa                 & lì-\textbf{t}óò             & li-\textbf{t}úβa                 & li-\textbf{t}ila                & í-\textbf{t}ù\textbf{t}ù               & mò-kù\textbf{t}à            &                        \\
KPE - SW                  & βé-\textbf{t}áa                 & lì-\textbf{t}óò             & li-\textbf{t}úβa                 & li-\textbf{t}ilá               & yú-\textbf{t}ù\textbf{t}ù              & mò-kù\textbf{t}à            & l-u\textbf{t}á                 \\
ISU - SW                  & βé-\textbf{t}ánu                & ì-\textbf{t}óì              & i-\textbf{t}úβa                  & i-\textbf{t}ila                 & yú-\textbf{t}ù\textbf{t}ù              & mò-kù\textbf{t}à            & i-ku\textbf{t}a                 \\
KOL - SW                  & bé-\textbf{t}ánu                & ì-\textbf{t}óì              & i-\textbf{t}úba                  & i-\textbf{t}ira                 & í-\textbf{t}ù\textbf{t}ù               & mò-kù\textbf{t}à            & i-wu\textbf{t}a                 \\
BUB - SW                  & βé-\textbf{t}áno                & lì-\textbf{t}óò             &                         & li-\textbf{t}ila                & í-\textbf{t}ù\textbf{t}ù               & mò-kù\textbf{t}à            & l-u\textbf{t}a                  \\
Kundu - N                 & be-\textbf{t}a                  & di-\textbf{t}oi                &                         & \textbf{t}ira                   & mo-\textbf{t}u\textbf{t}u                 &                       &        \\ 
   \lspbottomrule
 \end{tabular}
 }
\end{table}

\begin{table}[p]
\caption{Examples for the reflexes of PSB *k in C1 position}
\label{tab:grollemundbrownrule:kc1examples}
\fittable{
 \begin{tabular}{llllllll}
  \lsptoprule
\textbf{PSB}              & \textbf{*k$_1$}             & \textbf{*k$_1$}            & \textbf{*k$_1$}            & \textbf{*k$_1$ /*N\_}           & \textbf{*k$_1$ / \_*Vk$_2$/g$_2$}          & \textbf{*k$_1$ / *N\_}        & \textbf{*k$_1$ / \_*ek$_2$/g$_2$} \\
{{meaning}} & {{divide}} & {{storm}} & {{spear}} & {{tree trunk}} & {{sufficient/grow}} & {{pangolin}} & {{molar}}  \\
\midrule
BAN - S                   & i-\textbf{k}aɓa                   & ɓo-\textbf{h}uɗi                 & di-\textbf{h}ɔŋgɔ́               & ŋ-\textbf{k}ɔhɔ                       &                                   & \textbf{k}áŋgà                    & ɗi-\textbf{k}ɛhɔ                  \\
BAT - S                   & i-\textbf{k}aɓa                   & ɓo-\textbf{h}uɗi                 & di-\textbf{h}ɔŋgɔ́               & ŋ-\textbf{k}ɔhɔ                       & i-\textbf{k}owa                            & \textbf{k}áŋgà                    & ɗi-\textbf{k}ɛhɔ                  \\
YAS - S                   & e-\textbf{k}aɓa                   & ɓo-\textbf{k}uɗi                 & i-\textbf{k}ɔŋgɔ́                & mo-\textbf{k}ɔkɔ                      & e-\textbf{k}oka                            & \textbf{k}áŋgà                    & i-\textbf{c}ɔkɔ́                  \\
BAP - S                   & i-\textbf{k}aɓa                   & ɓo-\textbf{k}uɗi                 & i-\textbf{k}ɔŋgɔ́                & ŋ-\textbf{k}ɔkɔ                       &                                   & \textbf{k}áŋgà                    & i-\textbf{s}ɛkɔ                   \\
DUA - C                   & j-aɓa                    & m-udi                   & j-ɔŋgɔ́                 & mu-\textbf{k}ɔkɔ                      & \textbf{k}oka                              & \textbf{k}áàŋ                     & e-\textbf{k}ikɔ                   \\
BID - C                   & y-aɓa                    & m-uli                   & j-ɔŋgɔ́                 & ŋ-\textbf{k}ɔɣɔ                       & \textbf{k}oɣa                              & \textbf{k}áàŋ                     & kiɣi                     \\
OLI - C                   & y-aɓa                    & m-uli                   & j-ɔɣɔ́                  & ŋ-\textbf{k}ɔhɔ                       & \textbf{k}oɣa                              & \textbf{k}áàŋ                     & kiɣi                     \\
PON - C                   & ij-aɓa                   & m-udi                   & j-ɔŋgɔ́                 & mu-\textbf{k}ɔkɔ                      & i-\textbf{k}oka                            & \textbf{k}áàŋ                     & e-\textbf{k}ikɔ                   \\
MON - C                   & -\textbf{k}aɓanɛ                  & mu-\textbf{k}udi                 & j-ɔŋgɔ́                 & mu-\textbf{k}ɔkɔ                      & i-\textbf{k}oka                            & \textbf{k}áàŋ                     & e-\textbf{k}ikɔ                   \\
MUL - C                   & i-\textbf{k}aɓa                   & m-udi                   & j-ɔŋgɔ́                 & mu-\textbf{k}ɔɔ                       &                                   & \textbf{k}áàŋ                     & e-\textbf{c}ikɔ                   \\
MBO - SW                  &                          &                         & di-\textbf{k}ɔŋgɔ́               & mu-\textbf{k}ɔɔkɔ                     & li-\textbf{k}oka                           & \textbf{k}áŋgɛ̀                    & e-\textbf{k}ɛkɔ                   \\
KPE - SW                  &                          & mu-uli                  & li-ɔŋgɔ́                & mɔ-\textbf{k}ɔkɔ                      & li-\textbf{k}oka                           & \textbf{k}áŋgɛ̀                    & ɛ-\textbf{k}ɔkɔ                   \\
ISU - SW                  & i-\textbf{k}aβa                   & mu-\textbf{k}uli                 & di-\textbf{k}ɔŋgɔ́               & mo-\textbf{k}ɔkɔ                      & i-\textbf{k}oka                            & \textbf{k}áŋgɛ̀                    & e-\textbf{k}ikɔ                   \\
KOL - SW                  &                          &                         & j-ɔŋgɔ́                 & mo-\textbf{k}ɔkɔ                      & i-\textbf{k}oka                            & \textbf{k}áàŋ                     & e-\textbf{k}ikɔ                   \\
BUB - SW                  & li-\textbf{k}aβa                  &                         & li-ɔŋgɔ́                & mo-\textbf{k}ɔkɔ                      & li-\textbf{k}oka                           & \textbf{k}áŋgɛ̀                    & ɛ-\textbf{k}ɔkɔ                   \\
Kundu - N                 & -\textbf{k}aβa                    &                         & di-\textbf{k}ɔŋgɔ                & ɛ-\textbf{k}ɔkɔ (A11)                 &                                   &                            & e-\textbf{k}ekɔ         \\ 
   \lspbottomrule
 \end{tabular}
 }
\end{table}

\begin{table}[t]
\caption{Examples for the reflexes of PSB *k in C2 position}
\label{tab:grollemundbrownrule:kc2examples}
\fittable{
 \begin{tabular}{lllllllll}
  \lsptoprule
\textbf{PSB}              & \textbf{*k$_2$}               & \textbf{*k$_2$}          & \textbf{*k$_2$}            & \textbf{*k$_2$}          & \textbf{*k$_2$ / \_I}         & \textbf{*k$_2$}          & \textbf{*k$_2$}             & \textbf{*k$_2$}            \\
 {{meaning}} & {{elephant}} & {{one}} & {{scoop}} & {{fog}} & {{muteness}} & {{cow}} & {{mortar}} & {{spoon}} \\
\midrule
BAN - S                   & n-jɔ\textbf{h}u                     & y-ɔ\textbf{h}ɔ́                & i-tó\textbf{k}a                 & mbà\textbf{h}à               & ɗi-mbú\textbf{h}é                 & ɲa\textbf{k}a                  & è-ɓòwà                & tó\textbf{h}ò                  \\
BAT - S                   & n-jɔ\textbf{h}u                     & y-ɔ\textbf{h}ɔ́                & i-tó\textbf{k}a                 & mbà\textbf{h}à               & ɗi-mbú\textbf{h}é                 & ɲa\textbf{k}a                  & è-ɓòwà                & tó\textbf{h}ò                  \\
YAS - S                   & sɔ\textbf{k}u                       & e-vɔ́\textbf{k}ɔ́              &                         & mbà\textbf{k}à               & imbú\textbf{c}é                   & ɲa\textbf{k}a                  & è-ɓò\textbf{k}à                & tó\textbf{k}ò                  \\
BAP - S                   & n-jɔ\textbf{k}u                     & y-ɔ\textbf{k}ɔ́                & i-tó\textbf{k}a                 & mbà\textbf{k}à               & imbú\textbf{s}é                   & ɲa\textbf{k}a                  & è-ɓò\textbf{k}à                & tó\textbf{k}ò                  \\
DUA - C                   & n-jɔu                      & e-wɔ́                 & tóa                    & mbà                  & mbú\textbf{k}é                    & ɲa\textbf{k}a                  & è-ɓò\textbf{k}í                & tóò                   \\
BID - C                   & n-jɔ                       & fɔ́                   & tóa                    & mbà                  & mbú\textbf{ɣ}í                    & ɲa\textbf{ɣ}a                  & ɓò\textbf{ɣ}í                   & tóò                   \\
OLI - C                   & n-jɔ                       & wɔ́                   & tóa                    & mbà                  & mbú\textbf{ɣ}í                    & ɲa\textbf{ɣ}a                  & ɓò\textbf{ɣ}í                   & tóò                   \\
PON - C                   & n-jɔu                      & e-wɔ́                 & i-tóa                  & mbà                  & mbú\textbf{k}é                    & ɲa\textbf{k}a                  & è-ɓò\textbf{k}í                & tóò                   \\
MON - C                   & n-jɔ\textbf{k}u                     & y-ɔ\textbf{k}ɔ́                & i-tó\textbf{k}a                 & mbà\textbf{k}à               & mbú\textbf{k}é                    & ɲa\textbf{k}a                  & è-ɓò\textbf{k}í                & tó\textbf{k}ò                  \\
MUL - C                   & n-jɔu                      & e-wɔ́                 & i-tóa                  & mbà                  & mbú\textbf{c}é                    & ɲa\textbf{k}a                  & è-ɓò\textbf{k}í                & tóò                   \\
MBO - SW                  & n-jɔ\textbf{k}u                     & y-ɔ\textbf{k}ɔ́                & li-tó\textbf{k}a                & mbà\textbf{k}à               & mbú\textbf{k}é                    & ɲa\textbf{k}a                  & è-wò\textbf{k}í                &                         \\
KPE - SW                  & n-jɔ\textbf{k}u                     & y-ɔ\textbf{k}ɔ́                & li-tó\textbf{k}a                & mbà\textbf{k}ì               & mbú\textbf{k}é                    & ɲa\textbf{k}a                  & è-wò\textbf{k}í                & tó\textbf{k}ò                  \\
ISU - SW                  & n-jɔ\textbf{k}u                     & y-ɔ\textbf{k}ɔ́                & i-tó\textbf{k}a                 & mbà\textbf{k}à               & mbú\textbf{k}é                    & ɲa\textbf{k}a                  & è-ɓò\textbf{k}í                &                         \\
KOL - SW                  & n-jɔ\textbf{k}u                     & y-ɔ́\textbf{k}ɔ́               & i-tó\textbf{k}a                 & mbà\textbf{k}à               & mbú\textbf{k}é                    & ɲa\textbf{k}a                  & è-bò\textbf{k}í                & tó\textbf{k}ò                  \\
BUB - SW                  & n-jɔ\textbf{k}u                     & ey-ɔ́\textbf{k}ɔ́              & li-tó\textbf{k}a                & mbà\textbf{k}à               & mbú\textbf{k}é                    & ɲa\textbf{k}a                  &                          &                         \\
Kundu - N                 & n-jɔ\textbf{k}u                     & ey-ɔ\textbf{k}ɔ                &                         &                       &                            & ɲa\textbf{k}a                  &                          & to\textbf{k}o         \\ 
   \lspbottomrule
 \end{tabular}
 }
\end{table}

For the overlapping sets where we have k/c/s/Ø in C1 and C2 positions, we reconstructed one proto-sound *k based on several factors. First, some comparative sets can be explained by phonological contexts that triggered conditioned sound changes. For example, in C1 position, *k is maintained in Sawabantu languages when the stem presents another velar plosive (*k or *g) in C2 position (e.g., \tabref{tab:grollemundbrownrule:kc1examples}, examples `to be sufficient/to grow' or `molar') or when the words originally belonged to the PB class 9 \citep{hyman2019} with the nasal prefix (e.g., \tabref{tab:grollemundbrownrule:kc1examples} `pangolin', reconstructed *kákà in PB, BLR 3 1684, classes 9/10). When followed by a front vowel, *k$_{12}$ is palatalized in Yasa, Bapuku, Mulimba while the remaining languages present the regular reflexes for *k$_{12}$ (e.g., \tabref{tab:grollemundbrownrule:kc1examples} `molar' and \tabref {tab:grollemundbrownrule:kc2examples} `muteness'). Finally, we observe some sets for *k$_1$ with /h/ $\sim$ /k/ in Banoho and Batanga, /k/ in Yasa, Bapuku, Mongo, Mboko, Isu and Kundu, /Ø/ $\sim$ /k/ in Duala, Bidiman, Oli, Pongo, Mulimba, Kole, Mokpe and Bubia (e.g., \tabref{tab:grollemundbrownrule:kc1examples}, examples `divide', `storm', `spear'). The examination of the examples displaying these sound correspondences does not highlight a conditioning factor that would allow us to predict when we would have /k/, /h/ or Ø for the languages presenting two reflexes. Generally, in cases with similar partially overlapping correspondence sets with no phonological conditioning, the CM dictates that two proto-sounds should be reconstructed (e.g., *k$_1$ and *h$_1$). However, a closer examination of the data shows that for the languages with two reflexes, either /h/ $\sim$ /k/ or /Ø/ $\sim$ /k/, the sound sets with /k/ are attested in words that must have been acquired recently (e.g., `slave', `cassava') or words acquired by borrowings (e.g., `bag', `sword'). In sum, words that present the sound set with /k/ attested in all the languages are words that cannot be traced back to PB as opposed to the set with /h/ or /Ø/ (e.g., `hide', `root', `spear', `rope', etc.), also attested in a larger number of words. The study of the reflexes /h/ and /Ø/ highlights that PSB *k is undergoing a process of lenition in some Sawabantu languages until its complete deletion in most of the Central languages and some Southwest: PSB *k$_1$ > /h/ $\sim$ Ø. This type of lenition chain has been described by \citet[12]{pacchiarottiandbostoen2022} for West-Coastal Bantu languages spoken in southern Gabon, Congo and Democratic Republic of the Congo with Proto-West-Coastal Bantu *k$_2$ > x/ɣ/ʁ/h > Ø. In addition, \citet{grollemundandphilippsonforth2}, in a paper devoted to the study of the reflexes of PB *k and *g in NW Bantu languages (including A and some B languages), have observed that PB *k tends to be deleted with the same stages of lenition with generally PB *k$_{12}$ > /h/ > Ø. The lenition process is confirmed in C2 position with PSB *k$_2$ being lenited in the same languages (Banoho, Batanga, Duala, Bidiman, Oli, Pongo, Mulimba, Mokpe, examples are presented in \tabref{tab:grollemundbrownrule:kc2examples}). For Bidiman and Oli, we can observe that *k$_2$ > /ɣ/ (lenition) or Ø whereas in C1, the velar fricative does not appear (because PSB *k$_1$ has been deleted). In addition, we can notice that PSB *k$_{12}$ is always maintained and stable in Yasa, Bapuku, Mongo, Mboko, Isu and Kundu. Despite the multiple reflexes observed for PSB *k$_{12}$, our choice to reconstruct just one proto-sound was probably reasonable as we have observed that it was PSB *k that underwent lenition in some languages (while Yasa, Bapuku, Mongo, Mboko, Isu and Kundu only display /k/) and it is the same PSB *k that was retained in some possibly recently acquired words (that did not have the time to go through the lenition process yet). Finally, having two or more reflexes for a same proto-sound is a phenomenon attested in many Bantu languages with the double reflexes described by \citet{guthriecomparativebantu}, \citet{stewart1973, stewart1975}, \citet{janssens1991, janssens1992}, \citet{vanleynseeleandstewart1980} (see \citealt{philippson2022} for a detailed discussion on double reflexes found in NW Bantu languages) or multiple unconditioned reflexes (MUR) described by \citet{pacchiarottiandbostoen2022} referring to a situation where a same proto-sound yields two or more reflexes (without apparent conditioning).


We present in \tabref{tab:grollemundbrownrule:voicedplosivesexamples} the reflexes and proto-sounds reconstructed for voiced plosives. We added the reflexes for PSB *l$_{12}$ for analysis purposes. 

\begin{table}
\caption{Sound correspondences for voiced plosives and the approximant lateral in C1 and C2 positions. When two or more reflexes are presented, the first one is the most frequent. Bold characters mark sound changes from PSB to the daughter languages.}
\label{tab:grollemundbrownrule:voicedplosivesexamples}
\fittable{
 \begin{tabular}{lllllllll}
  \lsptoprule
\textbf{PSB} & \textbf{\begin{tabular}[c]{@{}l@{}}*b$_{12}$ / \\    \_i, u\end{tabular}} & \textbf{*b$_{12}$ / \_a,e,ɛ,ɔ,o} & \textbf{\begin{tabular}[c]{@{}l@{}}*b$_{12}$ / \\    \_u, o\end{tabular}} & \textbf{\begin{tabular}[c]{@{}l@{}}*b$_1$ /\\    \_wV\end{tabular}} & \textbf{*d$_1$} & \textbf{\begin{tabular}[c]{@{}l@{}}*d$_2$ / \\    \_i, u\end{tabular}} & \textbf{*d$_2$ /    \_a,e,ɛ,ɔ,o} & \textbf{*l$_{12}$} \\
\midrule
BAN - S      & \textbf{ɓ}                                                                    & \textbf{ɓ}                           & \textbf{ɓ}                                                                    & \textbf{ɓ(w)}                                                             & \textbf{ɗ}            & \textbf{ɗ}                                                                & \textbf{ɗ}                             & l             \\
BAT - S      & \textbf{ɓ}                                                                    & \textbf{ɓ}                           & \textbf{ɓ}                                                                    & \textbf{ɓ(w)}                                                             & \textbf{ɗ}            & \textbf{ɗ}                                                                & \textbf{ɗ}                             & l             \\
YAS - S      & \textbf{ɓ}                                                                    & \textbf{ɓ}                           & \textbf{ɓ}                                                                    & \textbf{ɓ(w)}                                                             & \textbf{ɗ}            & \textbf{ɗ}                                                                & \textbf{ɗ}                             & l             \\
BAP - S      & \textbf{ɓ}                                                                    & \textbf{ɓ}                           & \textbf{ɓ}                                                                    & \textbf{ɓ(w)}                                                             & \textbf{ɗ}            & \textbf{ɗ}                                                                & \textbf{ɗ}                             & l             \\
DUA - C      & b                                                                    & \textbf{ɓ}                           & b                                                                    & \textbf{ɓ(w)}                                                             & d            & d                                                                & \textbf{l}                             & l             \\
BID - C      & b                                                                    & \textbf{ɓ}                           & b                                                                    & \textbf{ɓ(w)}                                                             & d            & \textbf{l}                                                                & \textbf{l}                             & l             \\
OLI - C      & b                                                                    & \textbf{ɓ}                           & b                                                                    & \textbf{ɓ(w)}                                                             & d            & \textbf{l}                                                                & \textbf{l}                             & l             \\
PON - C      & b                                                                    & \textbf{ɓ}                           & b                                                                    & \textbf{ɓ(w)}                                                             & d            & d                                                                & \textbf{l}                             & l             \\
MON - C      & b                                                                    & \textbf{ɓ}                           & b                                                                    & \textbf{ɓ(w)}                                                             & d            & d                                                                & \textbf{l}                             & l             \\
MUL - C      & b                                                                    & \textbf{ɓ}                           & b                                                                    & \textbf{ɓ(w)}                                                             & d            & d                                                                & \textbf{ɗ}                             & l             \\
MBO - SW     & \textbf{β}                                                                    & \textbf{β}                           & \textbf{w $\sim$Ø}                                                            & \textbf{gb}                                                               & \textbf{l}            & \textbf{l}                                                                & \textbf{l}                             & l             \\
KPE - SW     & \textbf{β}                                                                    & \textbf{β}                           & \textbf{w $\sim$Ø}                                                            & \textbf{gb}                                                               & \textbf{l}            & \textbf{l}                                                                & \textbf{l}                             & \textbf{Ø} $\sim$l     \\
ISU - SW     & \textbf{β}                                                                    & \textbf{β}                           & \textbf{w $\sim$Ø}                                                            & \textbf{gb}                                                               & \textbf{l}            & \textbf{l}                                                                & \textbf{l}                             & l             \\
KOL - SW     & b                                                                    & b                           & b                                                                    & \textbf{gb}                                                               & \textbf{r}            & \textbf{r}                                                                & \textbf{r}                             & \textbf{r}             \\
BUB - SW     & \textbf{β}                                                                    & \textbf{β}                           & \textbf{w $\sim$Ø}                                                            & \textbf{gb}                                                               & \textbf{l}            & \textbf{l}                                                                & \textbf{l}                             & l $\sim$\textbf{Ø}     \\
Kundu - N    & b $\sim$\textbf{β}                                                            & b $\sim$\textbf{β}                   & b $\sim$\textbf{β}                                                            & \textbf{gb}                                                               & \textbf{r $\sim$l}    & \textbf{r $\sim$l}                                                        & \textbf{r $\sim$l}                     & \textbf{r}             \\ 
   \lspbottomrule
 \end{tabular}
 }
\end{table}

We reconstruct the proto-sound *b$_{12}$ for the partially overlapping comparative series ɓ/b/β and ɓ/b/w/β (e.g., \tabref{tab:grollemundbrownrule:voicedplosivesexamples}) based on the principle of directionality (lenition is more common than fortition, see above). We observe that PSB *b$_{12}$ > /ɓ/ in South languages, /b/ $\sim$ /ɓ/ in Central languages, /b/ or /β/ in North languages\footnote{\cite{friesen2002} often proposed the two variants with /b/ or /β/ for a same word.} and PSB *b$_{12}$ > the fricative /β/ in Southwest languages (lenition) except in Kole. The examples presented in \tabref{tab:grollemundbrownrule:bc1examples} and \tabref{tab:grollemundbrownrule:bc2examples} show that for Central languages, the distribution of /b/ and /ɓ/ is conditioned. The quality of the following vowel is going to trigger the sound to change: we will found the plosive /b/ before high vowels /\_i/ and /\_u/ (e.g., \tabref{tab:grollemundbrownrule:bc1examples} and \tabref{tab:grollemundbrownrule:bc2examples}) whereas the implosive /ɓ/ will occur elsewhere (/\_a, e, ɛ, ɔ, o/). This conditioning was already mentioned in \citet[18]{philippson2022} who demonstrated that PB *b does not present unconditioned double reflexes or MUR in Bantu A languages. For C1 and C2 in Southwest languages except for Kole, when followed by a back vowel (either \_o or \_u), PSB *b$_{12}$ will be deleted with sometimes an epenthetic glide /w/ added (e.g., \tabref{tab:grollemundbrownrule:bc1examples} and \tabref{tab:grollemundbrownrule:bc2examples}, examples `nine', `snot', `spider', `house', `hippo'). 

\begin{table}
\caption{Examples for the reflexes of PSB *b in C1 position}
\label{tab:grollemundbrownrule:bc1examples}
\fittable{
 \begin{tabular}{llllllll}
  \lsptoprule
\textbf{PSB}              & \textbf{*b$_1$}               & \textbf{*b$_1$}                                              & \textbf{*b$_1$}          & \textbf{*b$_1$}           & \textbf{*b$_1$\_ O, U}    & \textbf{*b$_1$ \_O, U}    & \textbf{*b$_1$ \_O, U}      \\
 {{meaning}} & {{marriage}} & {{two}}                                     & {{add}} & {{know}} & {{nine}} & {{snot}} & {{spider}} \\
\midrule
BAN - S                   & dì-\textbf{ɓ}áà                  & ɓé-\textbf{ɓ}à                                                   & ì-\textbf{ɓ}àɗɛ̀             & i-\textbf{ɓ}îa                 & ɗì-\textbf{ɓ}ùá              & ɗì-\textbf{ɓ}ɔ̀mbɔ́            & di-\textbf{ɓ}o\textbf{ɓ}ɛ                  \\
BAT - S                   & dì-\textbf{ɓ}áà                  & ɓé-\textbf{ɓ}à                                                   & ì-\textbf{ɓ}àɗɛ̀             & i-\textbf{ɓ}îa                 & ɗì-\textbf{ɓ}ùá              & ɗì-\textbf{ɓ}ɔ̀mbɔ́            & di-\textbf{ɓ}o\textbf{ɓ}ɛ                  \\
YAS - S                   & ì-\textbf{ɓ}áà                   & ɓí-\textbf{ɓ}à                                                   & è-\textbf{ɓ}àɗɛ̀             &                        & ɗì-\textbf{ɓ}ùá              & ì-\textbf{ɓ}ɔ̀mbɔ́             & di-\textbf{ɓ}o\textbf{ɓ}ɛ                  \\
BAP - S                   & ì-\textbf{ɓ}áà                   & ɓé-\textbf{ɓ}à                                                   & ì-\textbf{ɓ}àɗɛ̀             &                        & ɗì-\textbf{ɓ}ùá              & ɗì-\textbf{ɓ}ɔ̀mbɔ́            & di-\textbf{ɓ}o\textbf{ɓ}ɛ                  \\
DUA - C                   & dì-\textbf{ɓ}áà                  & ɓé-\textbf{ɓ}àá                                                 & \textbf{ɓ}átà                & \textbf{b}îa                    & dì-\textbf{b}ùá              & dì-\textbf{ɓ}ɔ̀mbɔ́            & di-\textbf{ɓ}o\textbf{ɓ}ɛ                  \\
BID - C                   & \textbf{b}áà                      & ɓí-\textbf{ɓ}àá                                                 & \textbf{ɓ}árɛ̀                & \textbf{b}îa                    & \textbf{b}ùa                   & \textbf{ɓ}ɔ̀mbɔ́                & di-\textbf{ɓ}o\textbf{ɓ}ɛ                  \\
OLI - C                   & \textbf{b}áà                      & ɓí-\textbf{ɓ}àá                                                 & \textbf{ɓ}árɛ̀                & \textbf{b}îa                    & \textbf{b}ùa                   & \textbf{ɓ}ɔ̀mbɔ́                & di-\textbf{ɓ}o\textbf{ɓ}ɛ                  \\
PON - C                   & ì-\textbf{ɓ}áà                   & ɓé-\textbf{ɓ}àá                                                 & ì-\textbf{ɓ}átà             & i-\textbf{b}îa                  & ì-\textbf{b}ùá               & ì-\textbf{ɓ}ɔ̀mbɔ́             & i-\textbf{ɓ}o\textbf{ɓ}ɛ                   \\
MON - C                   & dì-\textbf{ɓ}áà                  & ɓé-\textbf{ɓ}àá                                                 & ì-\textbf{ɓ}átà             & i-\textbf{b}îa                  & dì-\textbf{b}ùká             & dì-\textbf{ɓ}ɔ̀mbɔ́            & di-\textbf{ɓ}o\textbf{ɓ}ɛ                  \\
MUL - C                   & dì-\textbf{ɓ}áà                  & ɓé-\textbf{ɓ}àá                                                 & ì-\textbf{ɓ}átà             & i-\textbf{b}îa                  & dì-\textbf{b}ùá              & dì-\textbf{ɓ}ɔ̀mbɔ́            & di-\textbf{ɓ}o\textbf{ɓ}ɛ                  \\
MBO - SW                  & lì-\textbf{β}áà                  & βé-\textbf{β}àá                                                 & lì-\textbf{β}átà            &                        & lù-ùká              & \textbf{w}ɔ̀mbɔ́                & li-\textbf{w}o\textbf{β}ɛ                  \\
KPE - SW                  & lì-\textbf{β}áà                  & βé-\textbf{β}àá                                                 & lì-\textbf{β}átà            & li-\textbf{β}îa                & lù-ùwá              & lì-ɔ̀mbɔ́             & li-\textbf{w}o\textbf{β}ɛ                  \\
ISU - SW                  & lì-\textbf{β}áà                  & βé-\textbf{ɓ}àá                                                 & lì-\textbf{β}átà            & i-\textbf{b}îa                  & lì-\textbf{b}ùká             & lì-\textbf{w}ɔ̀mbɔ́            & li-\textbf{w}o\textbf{β}ɛ                  \\
KOL - SW                  & dì-\textbf{b}áà                  & bé-\textbf{b}àá                                                 & ì-\textbf{b}átà             & i-\textbf{b}îa                  & dì-\textbf{b}ùá              & dì-\textbf{b}ɔ̀mbɔ́            & di-\textbf{b}o\textbf{ɓ}ɛ                  \\
BUB - SW                  & lì-\textbf{β}áà                  & βé-\textbf{β}àá                                                 & lì-\textbf{β}átà            & li-\textbf{β}îa                & lì-gbá               & lì-ɔ̀mbɔ́             & li-\textbf{w}o\textbf{β}ɛ                  \\
Kundu - N                 &                            & \begin{tabular}[c]{@{}l@{}}be-\textbf{b}e/\\    βe-βe\end{tabular} &                       &                        &                        &                        & ri-\textbf{β}o\textbf{β}e              \\ 
   \lspbottomrule
 \end{tabular}
 }
\end{table}

\begin{table}
\caption{Examples for the reflexes of PSB *b in C2 position}
\label{tab:grollemundbrownrule:bc2examples}
\fittable{
 \begin{tabular}{lllllll}
  \lsptoprule
\textbf{PSB}              & \textbf{*b$_2$}          & \textbf{*b$_2$}             & \textbf{*b$_2$}             & \textbf{*b$_2$}                 & \textbf{*b$_2$ / \_o, u}   & \textbf{*b$_2$ / \_o, u}   \\
 {{meaning}} & {{six}} & {{divide}} & {{banana}} & {{wickedness}} & {{house}} & {{hippo}} \\
\midrule
BAN - S                   & n-tó\textbf{ɓ}á              & i-ka\textbf{ɓ}a                   & di-hú\textbf{ɓ}é                & ɓo\textbf{ɓ}é                        & ndá\textbf{ɓ}o                   & ŋɡu\textbf{ɓ}ú                  \\
BAT - S                   & n-tó\textbf{ɓ}á              & i-ka\textbf{ɓ}a                   & di-hú\textbf{ɓ}é                & ɓo\textbf{ɓ}é                        & ndá\textbf{ɓ}o                   & ŋɡu\textbf{ɓ}ú                  \\
YAS - S                   & mo-tóó\textbf{ɓ}á           & e-ka\textbf{ɓ}a                   & vi-kú\textbf{ɓ}é                &                              &                         & ŋɡu\textbf{ɓ}ú                  \\
BAP - S                   & n-tó\textbf{ɓ}á              & i-ka\textbf{ɓ}a                   & i-kú\textbf{ɓ}é                 & ɓo\textbf{ɓ}é                        & ndá\textbf{ɓ}o                   & ŋɡu\textbf{ɓ}ú                  \\
DUA - C                   & mu-tó\textbf{ɓ}á             & j-a\textbf{ɓ}a                    & di-kú\textbf{b}é                & ɓo\textbf{ɓ}é                        & ndá\textbf{ɓ}o                   & ŋɡu\textbf{b}ú                   \\
BID - C                   & n-tó\textbf{ɓ}á              & y-a\textbf{ɓ}a                    &                          & ɓo\textbf{b}î                        & ndá\textbf{ɓ}o                   & ŋɡu\textbf{b}ú                   \\
OLI - C                   & n-tó\textbf{ɓ}á              & y-a\textbf{ɓ}a                    &                          & ɓo\textbf{b}î                        & ndá\textbf{ɓ}o                   & ŋɡu\textbf{b}ú                   \\
PON - C                   & mu-tó\textbf{ɓ}á             & ij-a\textbf{ɓ}a                   & i-kú\textbf{ɓ}é                 & ɓo\textbf{ɓ}é                        & ndá\textbf{ɓ}o                   & ŋɡu\textbf{b}ú                   \\
MON - C                   & mu-tó\textbf{ɓ}á             & -ka\textbf{ɓ}anɛ                  & e-kú\textbf{ɓ}ékúɓé           & ɓo\textbf{ɓ}ée                       & ndá\textbf{ɓ}o                   & ŋɡu\textbf{b}ú                   \\
MUL - C                   & mu-tó\textbf{ɓ}á             & i-ka\textbf{ɓ}a                   & du-ú\textbf{ɓ}é                 & ɓo\textbf{ɓ}é                        & ndá\textbf{ɓ}o                   & ŋɡu\textbf{b}ú                   \\
MBO - SW                  & mo-tó\textbf{β}á             &                          &                          & wo\textbf{β}ée                       & ndá\textbf{w}o                   & ŋɡu\textbf{w}ú                  \\
KPE - SW                  & mo-tó\textbf{β}á             &                          &                          & wo\textbf{β}ée                       & ndá\textbf{w}o                   & ŋɡu\textbf{w}ú                  \\
ISU - SW                  & mo-tó\textbf{β}á             & i-ka\textbf{β}a                   &                          & wo\textbf{β}ée                       & ndá\textbf{w}o                   & ŋɡu\textbf{w}ú                  \\
KOL - SW                  & mo-tó\textbf{b}á             &                          &                          & bo\textbf{b}é                        & ndá\textbf{b}o                   & ŋɡu\textbf{b}ú                   \\
BUB - SW                  & mo-tó\textbf{β}á             & li-ka\textbf{β}a                  &                          & wo\textbf{β}ée                       & ndá\textbf{w}o                   & ŋɡu\textbf{w}ú                  \\
Kundu - N                 &                       & -ka\textbf{β}a                    &                          &                              & nda\textbf{b}o/nda\textbf{β}o             &                     \\ 
   \lspbottomrule
 \end{tabular}
 }
\end{table}

For the similar partially overlapping comparative sets with ɗ/d/l/r attested in C1 and C2 positions (e.g., \tabref{tab:grollemundbrownrule:dc1examples}), we chose to reconstruct two proto-sounds *d$_{12}$ and *l$_{12}$. The choice was complicated because we can observe the same sets of reflexes for Southwest languages where PSB *d$_{12}$ and *l$_{12}$ yield to /l/ (and /r/ in Kole). In addition, the study of the examples presented in \tabref{tab:grollemundbrownrule:dc1examples} and \tabref{tab:grollemundbrownrule:dc2examples} seems to display a phonological context that would require us to posit one proto-sound (PSB *d or *l) for these comparative sets with the impression that reflexes of *d$_{12}$ appear before high vowels (\_i and \_u) whereas the reflexes of *l$_{12}$ appear elsewhere. And finally, PB *d presenting two reflexes /d/\sim/l/ is a phenomenon that is attested in Bantu languages \citep{guthriecomparativebantu, meeussen1967, stewart1973, stewart1975, stewart1993, janssens1992, philippson2022}. 

\begin{table}
\caption{Examples for the reflexes of PSB *d and *l in C1 position}
\label{tab:grollemundbrownrule:dc1examples}
\fittable{
 \begin{tabular}{llllllll}
  \lsptoprule
\textbf{PSB}              & \textbf{*d$_1$}                 & \textbf{*d$_1$}          & \textbf{*d$_1$}            & \textbf{*d$_1$}          & \textbf{*l$_1$}           & \textbf{*l$_1$}           & \textbf{*l$_2$}           \\
 {{meaning}} & {{extinguish}} & {{old}} & {{water}} & {{run}} & {{send}} & {{read}} & {{three}} \\
\midrule
BAN - S                   & i-\textbf{ɗ}ima                       & ì-\textbf{ɗ}ùnà                & ma-\textbf{ɗ}îɓá                 & ì-\textbf{ɗ}àŋgwà           & ì-\textbf{l}óma               & ì-\textbf{l}áŋgà             & ɓé-\textbf{l}á\textbf{l}o               \\
BAT - S                   & i-\textbf{ɗ}ima                       & ì-\textbf{ɗ}ùnà                & ma-\textbf{ɗ}îɓá                 & ì-\textbf{ɗ}àŋgwà           & ì-\textbf{l}óma               & ì-\textbf{l}áŋgà             & ɓé-\textbf{l}á\textbf{l}o               \\
YAS - S                   & e-\textbf{ɗ}ima                       & è-\textbf{ɗ}ùnà                & me-héɓá                 &                       & è-\textbf{l}óma               & è-\textbf{l}áŋgà             & ɓî-\textbf{l}á\textbf{l}o                 \\
BAP - S                   & i-\textbf{ɗ}ima                       & ì-\textbf{ɗ}ùnà                & mi-iɓa                  &                       & ì-\textbf{l}óma               & ì-\textbf{l}áŋgà             & ɓé-\textbf{l}á\textbf{l}o               \\
DUA - C                   & \textbf{d}îma                         & \textbf{d}ùnà                  & ma-\textbf{d}îɓá                 & \textbf{ɗ}àŋgwà              & è-\textbf{l}óma               & \textbf{l}áŋgà                & ɓé-\textbf{l}á\textbf{l}o               \\
BID - C                   & \textbf{d}îmɛ                         & \textbf{d}ùnà                  & man-\textbf{d}îɓa                &                       & è-\textbf{l}óma               & \textbf{l}áŋgà                & bî-\textbf{l}á\textbf{l}o                 \\
OLI - C                   & \textbf{d}îmɛ                         & \textbf{d}ùnà                  & ma-\textbf{d}îɓá                 &                       & è-\textbf{l}óma               & \textbf{l}áɣà                 & bî-\textbf{l}á\textbf{l}o                 \\
PON - C                   & i-\textbf{d}îma                       & ì-\textbf{d}ùnà                & ma-\textbf{d}îɓá                 & ì-\textbf{ɗ}àŋgwà           & ì-\textbf{l}óma               & ì-\textbf{l}áŋgà             & ɓé-\textbf{l}á\textbf{l}o               \\
MON - C                   & i-\textbf{d}îma                       & ì-\textbf{d}ùnà                & ma-\textbf{d}îɓá                 &                       & ì-\textbf{l}óma               & ì-\textbf{l}áŋgà             & ɓé-\textbf{l}á\textbf{l}o               \\
MUL - C                   & i-\textbf{d}îma                       & ì-\textbf{d}ùnà                & ma-\textbf{d}îɓá                 & ì-\textbf{ɗ}àŋgwà           & ì-\textbf{l}óma               & ì-\textbf{l}áŋgà             & ɓé-\textbf{l}á\textbf{l}o               \\
MBO - SW                  & li-\textbf{l}îma                      & lù-ùnà                & ma-\textbf{d}îβá                 &                       & lì-\textbf{l}ómà             & lì-\textbf{l}áŋgà            & βé-\textbf{l}á\textbf{l}o                 \\
KPE - SW                  & li-\textbf{l}îma                      & lù-ùná                & ma-\textbf{d}îβá                 &                       & lò-óma               & lì-\textbf{l}áŋgà            & βé-\textbf{y}áo                  \\
ISU - SW                  & i-\textbf{l}îma                       & ì-\textbf{d}ùnà                & ma-\textbf{d}îβá                 &                       & lì-\textbf{l}ómà             & ì-\textbf{l}áŋgà             & βé-\textbf{l}á\textbf{l}o                 \\
KOL - SW                  & i-\textbf{d}îma                       & ì-\textbf{d}ùnà                & ma-\textbf{r}îβá                 &                       & ì-\textbf{r}óma               & ì-\textbf{r}áŋgà             & bé-\textbf{l}á\textbf{r}o               \\
BUB - SW                  & li-\textbf{l}îmɛ                      & li-\textbf{l}ùnà               & ma-\textbf{d}îβá                 &                       & lì-\textbf{y}óma              & lì-\textbf{y}áŋgà            & βé-\textbf{l}á\textbf{l}o                 \\
Kundu - N                 & \textbf{r}imɛ                         & unu                   & i-\textbf{l}iba/i-\textbf{r}iβa           &                       & \textbf{l}oma/\textbf{r}oma              &                        & be-\textbf{r}a\textbf{r}o            \\ 
   \lspbottomrule
 \end{tabular}
 }
\end{table}

\begin{table}
\caption{Examples for the reflexes of PSB *d and *l in C2 position}
\label{tab:grollemundbrownrule:dc2examples}
\fittable{
 \begin{tabular}{lllllllll}
  \lsptoprule
\textbf{PSB}              & \textbf{*d$_2$}             & \textbf{*d$_2$}            & \textbf{*d$_2$}           & \textbf{*d$_2$}                                              & \textbf{*d$_2$}                 & \textbf{*l$_2$}           & \textbf{*l$_2$}            & \textbf{*l$_2$}            \\
 {{meaning}} & {{turtle}} & {{birth}} & {{rope}} & {{goat}}                                    & {{put inside}} & {{lend}} & {{write}} & {{beard}}  \\
\midrule
BAN - S                   & kú\textbf{ɗ}ù                   & i-yá\textbf{ɗ}i                  & ŋ-hɔ\textbf{ɗ}î                &                        & ì-wɛ́\textbf{ɗ}ɛ̀                    & ì-ɓá\textbf{l}ɛ̀              & i-ti\textbf{l}a                  & ɗi-hɛ\textbf{l}ú                \\
BAT - S                   & kú\textbf{ɗ}ù                   & i-yá\textbf{ɗ}i                  & ŋ-hɔ\textbf{ɗ}î                & mbó\textbf{ɗ}ì                & ì-vɛ́\textbf{ɗ}ɛ̀                    & ì-ɓá\textbf{l}ɛ̀              & i-ti\textbf{l}a                  & ɗi-hɛ\textbf{l}ú                \\
YAS - S                   & kú\textbf{ɗ}ù                   & i-já\textbf{ɗ}i                  & mo-kɔ\textbf{ɗ}î               &                        & è-vɛ́\textbf{l}ɛ̀                    & è-ɓá\textbf{l}ɛ̀              & e-ti\textbf{l}a                  & sɛ\textbf{ɗ}u                    \\
BAP - S                   & kú\textbf{ɗ}ù                   & i-yá\textbf{ɗ}i                  & ŋ-kɔ\textbf{ɗ}î                &                        & ì-wɛ́\textbf{ɗ}ɛ̀                    & ì-ɓá\textbf{l}ɛ̀              & i-ti\textbf{l}a                  & i-hɛ\textbf{l}ú                 \\
DUA - C                   & wú\textbf{d}ù                   & yá\textbf{d}î                    & m-ɔ\textbf{d}î                 & mbó\textbf{d}ì                & wɛ́\textbf{l}ɛ̀                       & ɓá\textbf{l}ɛ̀                 & ti\textbf{l}a                    & sɛ\textbf{d}ú                   \\
BID - C                   & wú\textbf{l}ù                   & yá\textbf{l}é                    & m-ɔ\textbf{l}î                 & mbó\textbf{l}ì                & fɛ́\textbf{l}ɛ̀                       & ɓá\textbf{l}ɛ̀                 & ti\textbf{l}a                    & sɛ\textbf{l}ú                   \\
OLI - C                   & wú\textbf{l}ù                   & yá\textbf{l}î                    & m-ɔ\textbf{l}î                 & mbó\textbf{l}ì                & wɛ́\textbf{l}ɛ̀                       & ɓá\textbf{l}ɛ̀                 & ti\textbf{l}a                    & sɛ\textbf{l}ú                   \\
PON - C                   & wú\textbf{d}ù                   & ja\textbf{d}i                    & m-ɔ\textbf{d}î                 & mbó\textbf{d}ì                & ì-wɛ́\textbf{l}ɛ̀                    & ì-ɓá\textbf{l}ɛ̀              & i-ti\textbf{l}a                  & i-sɛ\textbf{d}ú                 \\
MON - C                   & kú\textbf{d}ù                   & ja\textbf{d}i                    & mu-kɔ\textbf{ɗ}î               & mbó\textbf{d}ì                & ì-wɛ́\textbf{l}ɛ̀                    & ì-ɓá\textbf{l}ɛ̀              & i-ti\textbf{l}a                  & i-sɛ\textbf{d}ú                 \\
MUL - C                   & kú\textbf{d}ù                   & i-yá\textbf{d}î                  & m-ɔ\textbf{d}î                 & mbó\textbf{d}ì                & ì-ɛ́\textbf{ɗ}ɛ̀                     & ì-ɓá\textbf{l}ɛ̀              & i-ti\textbf{l}a                  & di-hɛ\textbf{l}ú                \\
MBO - SW                  & è-kú\textbf{l}ɛ̀                & li-yá\textbf{l}î                 & mo-kɔ\textbf{l}î               & mbó\textbf{l}ì                & lì-wá\textbf{l}ɛ̀                   & lì-βá\textbf{l}ɛ̀             & li-ti\textbf{l}a                 & njɛ\textbf{l}i                   \\
KPE - SW                  & kú\textbf{l}ù                   & li-yá\textbf{l}î                 & mɔ-ɔ\textbf{l}î                & mbó\textbf{l}ì                &                              & lì-βá\textbf{l}ɛ̀             & li-ti\textbf{l}á                & njɛ\textbf{l}u                   \\
ISU - SW                  & è-kú\textbf{l}ɛ̀                & ja\textbf{l}i                    & mo-kɔ\textbf{l}î               & mbó\textbf{l}ì                &                              & ì-ɓá\textbf{l}ɛ̀              & i-ti\textbf{l}a                  & i-sɛ\textbf{l}ú                 \\
KOL - SW                  & è-kú\textbf{r}ɛ̀                & já\textbf{d}i                    & mo-kɔ\textbf{r}î               & mbó\textbf{r}ì                & ì-wɛ́\textbf{r}ɛ̀                    & ì-bá\textbf{r}ɛ̀              & i-ti\textbf{r}a                  & i-sɛ̀\textbf{r}ú                \\
BUB - SW                  & è-kú\textbf{l}ɛ̀                & li-yá\textbf{l}î                 & m-ɔ\textbf{l}î                 & mbó\textbf{l}ì                & lì-wá\textbf{l}ɛ̀                   & lì-ɓá\textbf{l}ɛ̀             & li-ti\textbf{l}a                 &                         \\
Kundu - N                 & ku                       & ja                      & mo-ko\textbf{r}i                & mbo\textbf{l}i/mbo\textbf{r}i            & u\textbf{l}e                          &                        & ti\textbf{r}a                    &              \\ 
   \lspbottomrule
 \end{tabular}
 }
\end{table}

In order to untangle this situation, we first need to summarize the reflexes attested for PSB *d$_{12}$ and *l$_{12}$ (see \tabref{tab:grollemundbrownrule:dc1examples} and \tabref{tab:grollemundbrownrule:dc2examples}). For South languages, we observe the reflexes /ɗ/ for *d$_{12}$ and /l/ for *l$_{12}$. For Central languages, PSB *d$_1$ > /d/ whereas in C2 position, we can find /d/ or /l/. In Duala, Pongo, Mongo, *d$_2$ > /d/ when followed by the high vowels / \_i, u whereas /l/ is attested elsewhere. For Mulimba, *d$_2$ > /d/ when followed by high-vowels and /ɗ/ is attested elsewhere. For Bidiman and Oli, *d$_2$ > /l/ in all contexts. For PSB *l$_{12}$, we found the reflexes /l/ for these languages. For Southwest languages, the situation is simpler with *d$_{12}$ > /l/ and *l$_{12}$ > /l/ in all the languages except for Kole that presents /r/ for both proto-sounds. We can note that Mokpe and Bubia can present the reflexes Ø or /l/ for PSB *l$_{12}$. Finally, in Kundu (North), the data presented in \cite[24]{friesen2002} gave several reflexes with *d$_{12}$ > /l/ or /r/ (\citealt[25-27]{friesen2002} explained that there is one phoneme /d/ in North languages that is realized [d] in front of the high vowels /i/ and /u/ whereas [l] is attested elsewhere with [r] occurring rarely). And *l$_{12}$ > /r/ in Kundu. 

\begin{table}[b]
\caption{Comparison of PSB *d$_2$ and *l$_2$.}
\label{tab:grollemundbrownrule:dandlc2examples}
\fittable{
 \begin{tabular}{lllllllll}
  \lsptoprule
\textbf{PSB}              & \textbf{*bodi}         & \textbf{*kado}           & \textbf{*bolo}         & \textbf{*pud}         & \textbf{*ul}              & \textbf{*vɛd}            & \textbf{*bal}          & \textbf{*ula}         \\
 {{meaning}} & {{goat}} & {{ladder}} & {{work}} & {{dig}} & {{be full}} & {{put in}} & {{lend}} & {{oil}} \\
\midrule
BAN - S                   &                        & ɓò-hà\textbf{ɗ}ò               &                        & i-pu\textbf{ɗ}a                & i᷅-hú\textbf{l}à                 & ì-wɛ́\textbf{ɗ}ɛ̀                & ì-ɓá\textbf{l}ɛ̀              & mà-ú\textbf{l}ɛ̀             \\
BAT - S                   & mbó\textbf{ɗ}ì                & ɓò-hà\textbf{ɗ}ò               &                        & i-pu\textbf{ɗ}a                & i᷅-hú\textbf{l}à                 & ì-vɛ́\textbf{ɗ}ɛ̀                & ì-ɓá\textbf{l}ɛ̀              & mà-hú\textbf{l}ɛ̀            \\
YAS - S                   &                        &                          & è-ɓó\textbf{l}ó              & e-pu\textbf{ɗ}wa               & è-hó\textbf{l}à                 & è-vɛ́\textbf{l}ɛ̀                & è-ɓá\textbf{l}ɛ̀              & màv-ú\textbf{l}ɛ̀            \\
BAP - S                   &                        & ɓò-kà\textbf{ɗ}ò               &                        & i-pu\textbf{ɗ}a                & i᷅-hú\textbf{l}à                 & ì-wɛ́\textbf{ɗ}ɛ̀                & ì-ɓá\textbf{l}ɛ̀              & má-\textbf{l}ɛ̀               \\
DUA - C                   & mbó\textbf{d}ì                &                          & è-ɓó\textbf{l}ó              & pu\textbf{l}a                  & jú\textbf{l}à                    & wɛ́\textbf{l}ɛ̀                   & ɓá\textbf{l}ɛ̀                 & mù-ú\textbf{l}à             \\
BID - C                   & mbó\textbf{l}ì                &                          & ɓò\textbf{l}ó                 &                       & ó\textbf{l}à                     & fɛ́\textbf{l}ɛ̀                   & ɓá\textbf{l}ɛ̀                 & mù-ú\textbf{l}à             \\
OLI - C                   & mbó\textbf{l}ì                &                          & ɓò\textbf{l}ó                 &                       & ó\textbf{l}à                     & wɛ́\textbf{l}ɛ̀                   & ɓá\textbf{l}ɛ̀                 & mù-ú\textbf{l}à             \\
PON - C                   & mbó\textbf{d}ì                &                          & è-ɓó\textbf{l}ó              & i-pu\textbf{l}a                & ì-jú\textbf{l}à                 & ì-wɛ́\textbf{l}ɛ̀                & ì-ɓá\textbf{l}ɛ̀              & mù-ú\textbf{l}à             \\
MON - C                   & mbó\textbf{d}ì                &                          & è-ɓó\textbf{l}ó              & i-pu\textbf{l}a                & ì-jú\textbf{l}à                 & ì-wɛ́\textbf{l}ɛ̀                & ì-ɓá\textbf{l}ɛ̀              & mà-wú\textbf{l}à            \\
MUL - C                   & mbó\textbf{d}ì                & ɓò-kà\textbf{ɗ}ò               & è-ɓó\textbf{l}ó              & i-pu\textbf{ɗ}a                & iú\textbf{l}à                    & ìɛ́\textbf{ɗ}ɛ̀                  & ì-ɓá\textbf{l}ɛ̀              & mù-ú\textbf{l}ɛ̀             \\
MBO - SW                  & mbó\textbf{l}ì                &                          & è-wò\textbf{l}ó              &                       & lù-ú\textbf{l}à                 & lì-wá\textbf{l}ɛ̀               & lì-βá\textbf{l}ɛ̀             & màl-úwà            \\
KPE - SW                  & mbó\textbf{l}ì                &                          & è-wò\textbf{l}ó              &                       & lù-ú\textbf{j}à                 &                          & lì-βá\textbf{l}ɛ̀             & mà-ú\textbf{j}à             \\
ISU - SW                  & mbó\textbf{l}ì                &                          & è-wò\textbf{l}ó              & i-ɸu\textbf{l}wa               & yù-ú\textbf{l}à                 &                          & ì-ɓá\textbf{l}ɛ̀              & mà-ú\textbf{l}à             \\
KOL - SW                  & mbó\textbf{r}ì                &                          & è-wò\textbf{r}ó              & i-ɸu\textbf{l}a                & i-wú\textbf{r}à                  & ì-wɛ́\textbf{r}ɛ̀                & ì-bá\textbf{r}ɛ̀              & mà-ú\textbf{r}à             \\
BUB - SW                  & mbó\textbf{l}ì                &                          & è-wò\textbf{l}ó              & li-ɸua                &                           & lì-wá\textbf{l}ɛ̀               & lì-ɓá\textbf{l}ɛ̀             &                       \\
Kundu - N                 & mbo\textbf{l}i                  &                          &                        &                       &                           & u\textbf{l}e                      &                        &                     \\ 
   \lspbottomrule
 \end{tabular}
 }
\end{table}

If we posit one proto-sound such as PSB *d$_{12}$ or *l$_{12}$ with the rule (conditioning) that the reflexes /d/ or /ɗ/ are attested in Sawabantu languages before high vowels whereas /l/ will be attested elsewhere, we would face several issues for the South and Central languages. First, for South languages, while we have observed that /ɗ/ is more often attested in front of high vowels whereas /l/ is more attested in front of non-high vowels, reconstructing one proto-sound would oblige us to posit double reflexes or MUR in C1 and C2 positions as /ɗ/ and /l/ can appear in similar phonological contexts (e.g., \tabref{tab:grollemundbrownrule:dc1examples} examples `run' and `read'; \tabref{tab:grollemundbrownrule:dandlc2examples}, examples `dig', `be full', `put in', `lend' or `oil'). Second, for Central languages, in C1 position, reconstructing one proto-sound is tempting as we found the reflexes /d/ only before high vowels and /l/ elsewhere. However, the situation in C2 with the conditioning described and more specifically the three reflexes observed in Mulimba with /d/, /ɗ/ and /l/ tipped the scales in favor of the reconstruction of two proto-sounds. In Mulimba, /d/ is found in front of high vowels (same as Duala, Pongo and Mongo) but /ɗ/ (as opposed to /l/ attested in Duala, Pongo and Mongo) is found in front of non-high vowels whereas /l/ is attested in many contexts, mostly in front of non-high vowels. The conditioning observed for /d/ and /ɗ/ along with the similar contexts in which we can find /ɗ/ and /l/ in Mulimba shows that these sets of reflexes should be linked to two separate proto-sounds PSB *d and *l. An examination of \tabref{tab:grollemundbrownrule:dandlc2examples} (along with \tabref{tab:grollemundbrownrule:dc1examples} and \tabref{tab:grollemundbrownrule:dc2examples}) confirms the need of two distinct proto-sounds which resulted in two sound sets: /d/ ~ /ɗ/ (< PSB *d) are attested in Mulimba when South languages also present /ɗ/ and reflexes /l/ (< PSB *l) are attested in Mulimba when South languages also present /l/.

Finally, there is some uncertainty as to the phonetic nature of the PSB (and PB) *b and *d as we cannot be absolutely sure that the original sounds were not implosives like in Southern and were later strengthened in some languages in front of first-degree vowels.



We do not reconstruct *g in PSB as the phoneme /g/ is absent in C1 and C2 positions in the daughter languages (except when preceded by a nasal homorganic prefix of class 9 such as in “moon” PSB *ŋ-gɔndɛ). This phenomenon is not unusual as it was noticed earlier by \citet{mohlig1981} and \citet{nurseandphilippson2003} who reported that Bantu languages spoken in the rainforest (generally A, B, C and parts of H) present this particular sound change where PB *g > /k/. In addition, \citet{pacchiarottiandbostoen2020} showed that for West-Coastal languages (parts of B and H), there was a velar merger that led PB *k and *g to merge into /k/. Finally, \citet{grollemundandphilippsonforth2} described that for NW Bantu languages (A and some B20), PB *k > Ø which triggered the PB *g > /k/. These sound changes modified the phonemic inventory of these Bantu languages yielding to the disappearance of /g/ in some languages. 

We also chose not to reconstruct labial-velar stops consonants *kp and *gb attested in Southwest and North Sawabantu languages (e.g., \tabref{tab:grollemundbrownrule:consonants} and \tabref{tab:grollemundbrownrule:voicedplosivesexamples}) as they are irregular and generally the result of contact or the fortition of the PSB sequences *kwV and PSB *bwV. Labial-velar stops are rare in Bantu languages but they do exist in several Bantu languages belonging to zones A, B and C. It is assumed that Bantu languages have acquired these labial-velar stops through contact with either Ubangian, Central Sudanic or Southern Bantoid languages \citep{clementsandrialland2008, bostoenanddonzo2013, idiatovandvandevelde2021}. Due the geographic location of Sawabantu languages in southern Cameroon, we could assume that Southwest and North Sawabantu languages have acquired these labial-velar stops by prolonged contact with Southern Bantoid languages located in the same region. While this hypothesis is highly plausible (and only a thorough study of these labial-velar stops will give us a clear answer), we have noticed some conditioning environments that might explain the presence of some labial-velar stops in Sawabantu languages. \citet[450]{bostoenanddonzo2013} examine the origin of labial-velar stops in Ngombe (C41), a Bantu language spoken in Democratic Republic of the Congo. While they do concur on the fact that labial-velar stops have been acquired by contact with neighboring Ubangian speakers, they also mention that conditioned sound changes can trigger labial-velar stops to develop. They demonstrate that labial-velars in Ngombe that do not have a Ubangian origin have developed because of the following rule: C$_{[+voiced]}$wV > gbV and C$_{[-voiced]}$wV > kpV, proposed first by \citet{Connell1994}, \citet{Connell1998} and \citet{cahill1999}. It seems that we have a similar phenomenon in Southwest and North Sawabantu languages where some labial-velar stops have possibly been acquired through contact while others seem to obey to this same conditioned sound change presented in \tabref{tab:grollemundbrownrule:labiovelars}. The status of these labial-velars stops /kp/ and /gb/ is still unclear and more data need to be analyzed. 

\begin{table}
\caption{Examples of the labial-velar stops in Southwest and North groups.}
\label{tab:grollemundbrownrule:labiovelars}
 \begin{tabular}{lllll}
  \lsptoprule
\textbf{PSB}       & \textbf{South} & \textbf{Central} & \textbf{Southwest} & \textbf{North}     \\
\midrule
*kwale “partridge” & /ŋgwae/        & /kwale/          & /\textbf{kp}ae/             & {{}} \\
*kwedi “death”     & /weɗi/         & /kweli/          & /\textbf{kp}eli/            &                    \\
*bwa “dog”         & /mbwa/         & /mbwa/           & /ŋ\textbf{gb}a/             & /m\textbf{gb}á/      \\ 
   \lspbottomrule
 \end{tabular}
\end{table}

For the fricatives presented in \tabref{tab:grollemundbrownrule:fricativesc1andc2}, we have reconstructed *v and *s in C1 and C2 positions. For the correspondence set with v/w/f/Ø, we chose to reconstruct *v since weakening (*v > /w/ > Ø) is more likely than strengthening. We could also have proposed *f but we hesitated to do so on the basis of Bidiman alone, which is rather idiosyncratic on other counts. For the set with h/s/ʒ, we reconstructed *s in C1 and C2 positions based on the same principle (with PSB *s$_{12}$ undergoing lenition to /h/ in South languages and also Mulimba). The exact phonetic value of the Mokpe sound /ʒ/ is the subject of some debate in the literature \citep[9]{atindogbe2013}.



\begin{table}
\caption{Sound correspondences for fricatives in C1 and C2 positions. Bold characters mark sound changes from PSB to the daughter languages.}
\label{tab:grollemundbrownrule:fricativesc1andc2}
 \begin{tabular}{lll}
  \lsptoprule
\textbf{PSB} & \textbf{*v$_{12}$} & \textbf{*s$_{12}$} \\
\midrule
BAN - S      & v             & \textbf{h}    \\
BAT - S      & v             & \textbf{h}    \\
YAS - S      & v             & \textbf{h}    \\
BAP - S      & v             & \textbf{h}    \\
DUA - C      & \textbf{w}    & s             \\
BID - C      & \textbf{f}    & s             \\
OLI - C      & v             & s             \\
PON - C      & \textbf{w}    & s             \\
   \lspbottomrule
\end{tabular}
\qquad
 \begin{tabular}{lll}
  \lsptoprule
\textbf{PSB} & \textbf{*v$_{12}$} & \textbf{*s$_{12}$} \\
\midrule
MON - C      & \textbf{w}    & s             \\
MUL - C      & \textbf{Ø}    & \textbf{h}    \\
MBO - SW     & \textbf{Ø}    & s             \\
KPE - SW     & \textbf{Ø}    & \textbf{ʒ}    \\
ISU - SW     & \textbf{Ø}    & s             \\
KOL - SW     & \textbf{w}    & s             \\
BUB - SW     & \textbf{w}    & s             \\
Kundu - N    & \textbf{Ø}    & s       \\ 
   \lspbottomrule
 \end{tabular}
\end{table}

\begin{table}
\caption{Examples for the reflexes of PSB *v and *s in C1 and C2 positions.}
\label{tab:grollemundbrownrule:vandsc1andc2}
\fittable{
 \begin{tabular}{lllllllll}
  \lsptoprule
\textbf{PSB}              & \textbf{*v$_1$}           & \textbf{*v$_1$}           & \textbf{*v$_2$}           & \textbf{*s$_1$}           & \textbf{*s$_1$}           & \textbf{*s$_2$}           & \textbf{*s$_2$}             \\
 {{meaning}} & {{wind}} & {{twin}} & {{itch}} & {{iron}} & {{army}} & {{corn}} & {{choose}} \\
\midrule
BAN - S                   & ɗi-\textbf{v}îndî               & ɗì-\textbf{v}áhà             & yò\textbf{v}á                 & i-\textbf{h}îa                  & \textbf{h}áo                   & mbá\textbf{y}í                &                          \\
BAT - S                   & ɗi-\textbf{v}îndî               & ɗì-\textbf{v}áhà             & yò\textbf{v}á                 & i-\textbf{h}îa                  &                        & mbá\textbf{h}í                & i-pɔ\textbf{h}ɔ                   \\
YAS - S                   & i-\textbf{v}îndó                & ì-\textbf{v}ágà              & ò\textbf{v}á                  & e-\textbf{h}îsa                 & i-\textbf{y}áko                & mbà\textbf{s}ì                &                          \\
BAP - S                   & ɗi-\textbf{v}îndî               & ì-\textbf{v}áhà              & yò\textbf{v}á                 & i-\textbf{h}îa                  & \textbf{h}áo                   & mbá\textbf{y}í                &                          \\
DUA - C                   & di-\textbf{w}îndî               & dì-\textbf{w}ásà             &                        & \textbf{s}îa                    & \textbf{s}áo                   & mbà\textbf{s}ì                & pɔ\textbf{s}ɔ                     \\
BID - C                   & \textbf{f}indîfindî             & \textbf{f}ásà                 &                        & \textbf{s}îa                    & \textbf{s}áo                   & mbà\textbf{s}ì                & pɔ\textbf{s}ɔ                     \\
OLI - C                   & \textbf{v}îndîvîndî             & \textbf{v}ásà                 &                        & \textbf{s}îa                    & \textbf{s}áo                   & mbà\textbf{s}ì                & pɔ\textbf{s}ɔ                     \\
PON - C                   & i-\textbf{w}îndîwîndî           & ì-\textbf{w}ásà              &                        & i-\textbf{s}îa                  & \textbf{s}áo                   & mbà\textbf{s}ì                & i-pɔ\textbf{s}ɔ                   \\
MON - C                   & di-\textbf{w}îndîwîndî          & dì-\textbf{w}ásà             &                        & i-\textbf{s}îa                  & \textbf{s}áko                  & mbà\textbf{s}ì                & i-pɔ\textbf{s}ɔ                   \\
MUL - C                   & di-îndî                & dì-áhà              & yò\textbf{w}á                 & i-\textbf{h}îa                  & \textbf{h}áo                   & mbà\textbf{s}ì                & i-pɔ\textbf{h}ɔ                   \\
MBO - SW                  & li-\textbf{w}îndó               & lì-ásɛ̀              &                        & li-\textbf{s}îa                 & \textbf{s}áko                  &                        & li-ɸɔ\textbf{s}ɔ́                 \\
KPE - SW                  & li-\textbf{w}îndî               & mà-ɛ́ʒɛ̀              &                        & li-\textbf{ʒ}îa                 & \textbf{ʒ}iyaó                 & mbà\textbf{ʒ}ì                & li-ɸɔ\textbf{ʒ}ɔ́                 \\
ISU - SW                  & li-\textbf{w}îndîwîndî          & mà-ɛ́sɛ̀              &                        & i-\textbf{s}îa                  & \textbf{s}áo                   & mbà\textbf{s}ì                & li-ɸɔ\textbf{s}ɔ                  \\
KOL - SW                  & di-\textbf{w}îndî               & dì-\textbf{w}ásà             &                        & i-\textbf{s}îa                  & \textbf{s}íáko                & mbà\textbf{s}ì                & li-ɸɔ\textbf{s}ɔ                  \\
BUB - SW                  & li-îndó                & m-ásɛ̀                &                        & li-\textbf{s}îa                 & \textbf{s}íyaó                & mbà\textbf{s}ì                & li-ɸɔ\textbf{s}ɔ                  \\
Kundu - N                 &                        & ri-asɛ                 &                        &                        &                        &                        & βɔ\textbf{s}ɔ                      \\ 
   \lspbottomrule
 \end{tabular}
 }
\end{table}

Nasal consonants *m, *n, *ɲ have been reconstructed in C1 and C2 position along with prenasalized consonants *mb, *nd, *ŋg (e.g., \tabref{tab:grollemundbrownrule:nasalscorrespondence}). 


\begin{table}[p]

\caption{Sound correspondences for nasals and prenasalized consonants. Bold characters mark sound changes from PSB to the daughter languages. We can note that denasalization is frequent in Oli.}
\label{tab:grollemundbrownrule:nasalscorrespondence}
 \begin{tabular}{lllllllll}
  \lsptoprule
\textbf{PSB} & \textbf{*m$_{12}$} & \textbf{*n$_{12}$} & \textbf{*ɲ$_{12}$} & \textbf{*mb} & \textbf{*nd} & \textbf{*ŋg} \\
\midrule
BAN - S      & m             & n             & ɲ             & mb           & nd           & ŋg           \\
BAT - S      & m             & n             & ɲ             & mb           & nd           & ŋg           \\
YAS - S      & m             & n             & ɲ             & mb           & nd           & ŋg           \\
BAP - S      & m             & n             & ɲ             & mb           & nd           & ŋg           \\
DUA - C      & m             & n             & ɲ             & mb           & nd           & ŋg           \\
BID - C      & m             & n             & ɲ             & mb           & nd           & ŋg           \\
OLI - C      & m             & n             & ɲ             & mb           & nd           & \textbf{ɣ} $\sim$ŋg   \\
PON - C      & m             & n             & ɲ             & mb           & nd           & ŋg           \\
MON - C      & m             & n             & ɲ             & mb           & nd           & ŋg           \\
MUL - C      & m             & n             & ɲ             & mb           & nd           & ŋg           \\
MBO - SW     & m             & n             & ɲ             & mb           & nd           & ŋg           \\
KPE - SW     & m             & n             & ɲ             & mb           & nd           & ŋg           \\
ISU - SW     & m             & n             & ɲ             & mb           & nd           & ŋg           \\
KOL - SW     & m             & n             & ɲ             & mb           & nd           & ŋg           \\
BUB - SW     & m             & n             & ɲ             & mb           & nd           & ŋg           \\
Kundu - N    & m             & n             & ɲ             & mb           & nd           & ŋg           \\
   \lspbottomrule
 \end{tabular}
\end{table}

\begin{table}[p]

\caption{Examples for the reflexes of PSB *m, *n and *ɲ in C1 and C2 positions.}
\label{tab:grollemundbrownrule:nasalexamples}
\fittable{
 \begin{tabular}{lllllllll}
  \lsptoprule
\textbf{PSB}              & \textbf{*m$_1$}             & \textbf{*m$_2$}            & \textbf{*n$_{12}$}          & \textbf{*n$_2$}            & \textbf{*ɲ$_1$}                & \textbf{*ɲ$_1$}                & \textbf{*ɲ$_2$}            \\
 {{meaning}} & {{finish}} & {{stand}} & {{bird}} & {{shame}} & {{surprised}} & {{happiness}} & {{shine}} \\
\midrule
BAN - S                   & ì-\textbf{m}à                     & i-tɛ́\textbf{m}ɛ                 & i-\textbf{n}ɔ\textbf{n}í                & ì-yɔ́\textbf{n}ì               &                             & e-\textbf{ɲ}ɛŋɡɛ                     & ì-pà\textbf{ɲ}à               \\
BAT - S                   & ì-\textbf{m}à                     & i-tɛ́\textbf{m}ɛ                 & i-\textbf{n}ɔ\textbf{n}í                & ì-hɔ́\textbf{n}ì               &                             & e-\textbf{ɲ}ɛŋɡɛ                     & ì-pà\textbf{ɲ}à               \\
YAS - S                   &                          & e-tɛ́\textbf{m}ɛ                 & vi-ɲɔ\textbf{n}í               & vì-yɔ́\textbf{n}ì              &                             & e-\textbf{ɲ}ɛŋɡɛ                     & è-pà\textbf{ɲ}à               \\
BAP - S                   & ì-\textbf{m}à                     & i-tɛ́\textbf{m}ɛ                 & i-\textbf{n}ɔ\textbf{n}í                & ì-yɔ́\textbf{n}ì               &                             & e-\textbf{ɲ}ɛŋɡɛ                     & ì-pà\textbf{ɲ}à               \\
DUA - C                   &                          & tɛ́\textbf{m}ɛ                   & i-\textbf{n}ɔɔ́\textbf{n}                & ì-sɔ́ɔ̀\textbf{n}               & \textbf{ɲ}áka                        & mu-\textbf{ɲ}ɛŋɡɛ                    & pà\textbf{ɲ}a                   \\
BID - C                   & \textbf{m}à                       & tɛ́\textbf{m}ɛ                   & \textbf{n}ɔɔ́\textbf{n}                  & sɔ́ɔ̀\textbf{n}                  & \textbf{ɲ}áɣa                       & ɲ̀-\textbf{ɲ}ɛŋɡɛ                    & pà                     \\
OLI - C                   & \textbf{m}à                       & tɛ́\textbf{m}ɛ                   & \textbf{n}ɔɔ́\textbf{n}                  & sɔ́ɔ̀\textbf{n}                  & \textbf{ɲ}áɣa                       & ɲ̀-\textbf{ɲ}ɛŋɡɛ                    & pà                     \\
PON - C                   &                          & i-tɛ́\textbf{m}ɛ                 & e-\textbf{n}ɔɔ́\textbf{n}                & è-sɔ́ɔ̀\textbf{n}               & i-\textbf{ɲ}áka                      & mu-\textbf{ɲ}ɛŋɡɛ                    & ì-pà\textbf{ɲ}à               \\
MON - C                   & ì-\textbf{m}à                     & i-tɛ́\textbf{m}ɛ                 & i-\textbf{n}ɔɔ́\textbf{n}                & ì-sɔ́ɔ̀\textbf{n}               & i-\textbf{ɲ}áka                      & mu-\textbf{ɲ}ɛŋɡɛ                    & ì-pà\textbf{ɲ}à               \\
MUL - C                   &                          & i-tɛ́\textbf{m}ɛ                 & i-\textbf{n}ɔɔ́\textbf{n}                & ì-hɔ́ɔ̀\textbf{n}               & i-\textbf{ɲ}áka                      & mu-\textbf{ɲ}ɛnjɛ                    & ì-pà\textbf{ɲ}à               \\
MBO - SW                  &                          & li-tɛ́\textbf{m}ɛ                & i-\textbf{n}ɔ\textbf{n}í                & ì-sɔ́\textbf{n}ì               & li-\textbf{ɲ}áka                     & mo-\textbf{ɲ}ɛŋɡi                    & lì-kpà\textbf{ɲ}à             \\
KPE - SW                  &                          & li-tɛ́\textbf{m}ɛ                & i-\textbf{n}ɔ\textbf{n}í                & ì-ʒɔ́\textbf{n}ì               & li-\textbf{ɲ}áka                     & mo-\textbf{ɲ}ɛŋɡi                    &                         \\
ISU - SW                  & ì-\textbf{m}à                     & i-tɛ́\textbf{m}ɛ                 & i-\textbf{n}ɔ\textbf{n}í                & ì-sɔ́ɔ̀\textbf{n}               & i-\textbf{ɲ}áka                      & mo-\textbf{ɲ}ɛŋɡi                    & ì-pà\textbf{ɲ}à               \\
KOL - SW                  & ì-\textbf{m}à                     & i-tɛ́\textbf{m}ɛ                 & i-\textbf{n}ɔ\textbf{n}í                & ì-sɔ́\textbf{n}ì               & i-\textbf{ɲ}áka                      & mo-\textbf{ɲ}ɛŋɡi                    & ì-pà\textbf{ɲ}à               \\
BUB - SW                  &                          & li-tɛ́\textbf{m}ɛ                & i-\textbf{n}ɔ\textbf{n}í                & ì-sɔ́\textbf{n}ì               & li-\textbf{ɲ}áka                     & mo-\textbf{ɲ}ɛŋɡi                    &                         \\
Kundu - N                 & \textbf{m}arɛ                     &                         & i-\textbf{n}ɔ                   & i-sɔ                    & ma-\textbf{ɲ}aka                     &                             &               \\ 
   \lspbottomrule
 \end{tabular}
 }
\end{table}

\clearpage
Prenasalized consonants can also be found in C1 position as they are the result of the combination of the nasal prefix of class 9 *N- (homorganic nasal) combined with a plosive consonant. For example, the word for `crocodile' is reconstructed *gàndʊ́ in PB (BLR3 1326) with a class 9 which led to the PB root *N-gàndʊ́ > [ŋ-gàndó] in Banoho. 

\begin{table}
\caption{Examples for the reflexes of PSB *mb, *nd and *ŋg.}
\label{tab:grollemundbrownrule:prenasals}
 \begin{tabular}{llll}
  \lsptoprule
\textbf{PSB}              & \textbf{*mb}               & \textbf{*nd}                & \textbf{*ŋg}           \\
 {{meaning}} & {{namesake}} & {{crocodile}} & {{root}} \\
\midrule
BAN - S                   & \textbf{mb}ó\textbf{mb}ò                   & ŋgà\textbf{nd}ó                    & ŋ̀-hà\textbf{ŋg}á             \\
BAT - S                   & \textbf{mb}ó\textbf{mb}ò                   & ŋgà\textbf{nd}ó                    & ŋ̀-hà\textbf{ŋg}á             \\
YAS - S                   & \textbf{mb}ó\textbf{mb}ò                   & ŋgà\textbf{nd}ó                    &                        \\
BAP - S                   & \textbf{mb}ó\textbf{mb}ò                   & ŋgà\textbf{nd}ó                    & ŋ̀-kà\textbf{ŋg}á             \\
DUA - C                   & \textbf{mb}ó\textbf{mb}ò                   & ŋgà\textbf{nd}ó                    & mw-à\textbf{ŋg}á              \\
BID - C                   & \textbf{mb}ó\textbf{mb}ò                   & ŋgà\textbf{nd}ó                    & mw-à\textbf{ŋg}á              \\
OLI - C                   & \textbf{mb}ó\textbf{mb}ò                   & ŋgà\textbf{nd}ó                    & mw-à\textbf{ɣ}á               \\
PON - C                   & \textbf{mb}ó\textbf{mb}ò                   & ŋgà\textbf{nd}ó                    & mw-à\textbf{ŋg}á              \\
MON - C                   & \textbf{mb}ó\textbf{mb}ò                   & ŋgà\textbf{nd}ó                    & mù-kà\textbf{ŋg}á            \\
MUL - C                   & \textbf{mb}ó\textbf{mb}ò                   & ŋgà\textbf{nd}ó                    & mw-à\textbf{ŋg}á              \\
MBO - SW                  & \textbf{mb}ó\textbf{mb}ò                   & ŋgà\textbf{nd}ó                    & ŋm-à\textbf{ŋg}á              \\
KPE - SW                  & \textbf{mb}ó\textbf{mb}ò                   & ŋgà\textbf{nd}ó                    & ŋm-à\textbf{ŋg}á              \\
ISU - SW                  & \textbf{mb}ó\textbf{mb}ò                   & ŋgà\textbf{nd}ó                    & mò-kà\textbf{ŋg}á            \\
KOL - SW                  & \textbf{mb}ó\textbf{mb}ò                   & ŋgà\textbf{nd}ó                    & ŋw-à\textbf{ŋg}á              \\
BUB - SW                  & \textbf{mb}ó\textbf{mb}ò                   & ŋgà\textbf{nd}ó                    & ŋm-à\textbf{ŋg}á              \\
Kundu - N                 &                            & ŋga\textbf{nd}o                      &     \\ 
   \lspbottomrule
 \end{tabular}
\end{table}

We also reconstructed the approximants *w and *y (presented in \tabref{tab:grollemundbrownrule:approximants}) in C1 and C2 positions. 


\begin{table}
\caption{Sound correspondences for approximants *w and *y.}
\label{tab:grollemundbrownrule:approximants}
 \begin{tabular}{lll}
  \lsptoprule
\textbf{PSB} & \textbf{*w} & \textbf{*y} \\
\midrule
BAN - S      & w           & y           \\
BAT - S      & w           & y           \\
YAS - S      & w           & y           \\
BAP - S      & w           & y           \\
DUA - C      & w           & y           \\
BID - C      & w           & y           \\
OLI - C      & w           & y           \\
PON - C      & w           & y           \\
   \lspbottomrule
 \end{tabular}
 \qquad
 \begin{tabular}{lll}
  \lsptoprule
\textbf{PSB} & \textbf{*w} & \textbf{*y} \\
\midrule
MON - C      & w           & y           \\
MUL - C      & w           & y           \\
MBO - SW     & w           & y           \\
KPE - SW     & w           & y           \\
ISU - SW     & w           & y           \\
KOL - SW     & w           & y           \\
BUB - SW     & w           & y           \\
Kundu - N    & w           & y           \\
   \lspbottomrule
 \end{tabular}
\end{table}

The study of *w$_{12}$ and *y$_{12}$ shows that in many examples, /w/ and /y/ are epenthetic consonants in Sawabantu languages resulting from various processes occurring at morpheme boundaries. Generally, when PSB *k deletes, it leads to the addition of /w/ or /y/ (depending on the quality of the following vowel) in between the prefix and the root: PSB *kuba `hen' > /wuɓa/ in Duala, Bidiman, Oli and Pongo. Also, when PSB *b > Ø / \_o, u, it leads to the addition of /w/: PSB *bɔmbɔ `snot' > /wɔmbɔ/ in Mbole. Finally, we also have examples where PSB *v > /w/ in some languages: PSB *vasa `twin' > /wasa/ in Pongo. However, there are also cases where *w can be reconstructed in C1 and C2 position (e.g., \tabref{tab:grollemundbrownrule:approximantexamples}) or in between C1 and V1 (e.g., PSB *bwa `dog' > /imbwa/ in Batanga for example) or C2 and V2 (e.g., PSB *puŋgwa `mix' > /i-puŋgwa/ in Banoko for example). Finally, we needed to reconstruct *y$_1$ for comparative series where /y/ is attested in all Sawabantu languages (e.g., \tabref{tab:grollemundbrownrule:approximantexamples}, examples PSB *ya `give birth' or PSB *yoko `penis'). We did not find non-epenthetic *y in C2 position. 

\begin{table}
\caption{Examples for the reflexes of PSB *w and *y in C1 and C2 positions.}
\label{tab:grollemundbrownrule:approximantexamples}
 \begin{tabular}{lllll}
  \lsptoprule
\textbf{PSB}              & \textbf{*w$_1$}          & \textbf{*w$_2$}            & \textbf{*y$_1$}            & \textbf{*y$_1$}                 \\
 {{meaning}} & {{die}} & {{coast}} & {{penis}} & {{give birth}} \\
\midrule
BAN - S                   & i-\textbf{w}ɔ́ɔ                &                         & \textbf{y}òó                   & i-\textbf{y}áa                        \\
BAT - S                   & i-\textbf{w}áa                &                         & \textbf{y}òó                   & i-\textbf{y}áa                        \\
YAS - S                   & e-\textbf{w}áa                &                         &                         & e-\textbf{j}áa                        \\
BAP - S                   & i-\textbf{w}áa                &                         & \textbf{y}òó                   & i-\textbf{y}áa                        \\
DUA - C                   & \textbf{w}ɔ́                   & sá\textbf{w}á                  & è-\textbf{y}òó                & \textbf{y}áa                          \\
BID - C                   & \textbf{w}áa                  & sá\textbf{w}á                  & \textbf{y}òó                   & \textbf{y}áa                          \\
OLI - C                   & \textbf{w}áa                  & sá\textbf{w}á                  & \textbf{y}òó                   & \textbf{y}áa                          \\
PON - C                   & ij-ɔ́ɔ                & sá\textbf{w}á                  & è-\textbf{y}òó                & i-\textbf{y}áa                        \\
MON - C                   & i-\textbf{w}áa                & sá\textbf{w}á                  & è-\textbf{y}òkó               & i-\textbf{y}áa                        \\
MUL - C                   & i-\textbf{w}ɔ́                 & sá\textbf{w}á                  & è-\textbf{y}òó                & i-\textbf{y}áa                        \\
MBO - SW                  & li-\textbf{w}áa               & sá\textbf{w}á                  & \textbf{y}òkó                  & li-\textbf{j}áa                       \\
KPE - SW                  & liŋm-áa              & ʒá\textbf{w}á                  & \textbf{y}òkó                  & li-\textbf{y}áa                       \\
ISU - SW                  & i-\textbf{w}áa                & sá\textbf{w}á                  & \textbf{y}òkó                  & i-\textbf{y}áa                        \\
KOL - SW                  & i-\textbf{w}áa                & sá\textbf{w}á                  & è-\textbf{y}òkó               & i-\textbf{y}áa                        \\
BUB - SW                  & li-\textbf{w}áa               & sá\textbf{w}á                  & è-\textbf{y}òkó               & li-\textbf{y}ɛ́ɛ                      \\
Kundu - N                 & \textbf{w}a                    &                         &                         &   \\
   \lspbottomrule
 \end{tabular}
\end{table}

We do find in Ebobissé’s data (\citeyear{ebobisse2015}) the occurrence of /j/ (/ʤ/) that appears rarely and for this reason, we have decided to not reconstruct *j as it is often a prefix attached to a root that starts with a vowel (e.g., \tabref{tab:grollemundbrownrule:jexamples} `ten' or `mirror') or in free variation with /y/ for which we have reconstructed *y (e.g., \tabref{tab:grollemundbrownrule:jexamples} “wash”).

\begin{table}
\caption{Examples with /j/ in Sawabantu languages.}
\label{tab:grollemundbrownrule:jexamples}
 \begin{tabular}{llll}
  \lsptoprule
\textbf{PSB}              & \textbf{-}            & \textbf{-}               & \textbf{*y$_1$}           \\
 {{meaning}} & {{ten}} & {{mirror}} & {{wash}} \\
\midrule
BAN - S                   & \textbf{j}óòm                & \textbf{j}ɛní                    & i-\textbf{y}oa                  \\
BAT - S                   & \textbf{j}óòm                & \textbf{j}ɛní                    & i-\textbf{y}oa                  \\
YAS - S                   & \textbf{j}óòm                & \textbf{j}ɛní                    &                        \\
BAP - S                   & \textbf{j}óòm                & \textbf{j}ɛní                    & i-\textbf{y}oa                  \\
DUA - C                   & ɗ-óòm               & \textbf{j}ɛnɛ́                    & \textbf{j}óa                   \\
BID - C                   & \textbf{j}óòm                & \textbf{j}ɛnɛ́                    & \textbf{y}oa                    \\
OLI - C                   & \textbf{j}óòm                & \textbf{j}ɛnɛ́                    & \textbf{y}oa                    \\
PON - C                   & ɗ-óòm               & \textbf{j}ɛnɛ́                    & i-\textbf{j}oa                  \\
MON - C                   &                       & \textbf{j}ɛnɛ́                    & i-\textbf{j}óa                 \\
MUL - C                   & \textbf{j}óòm                & \textbf{j}ɛnɛ́                    & i-\textbf{y}oa                  \\
MBO - SW                  & dì-ómè             &                          &                        \\
KPE - SW                  & lì-\textbf{y}ómè            & li-ɛní                  &                        \\
ISU - SW                  &                       & \textbf{j}ɛní                    & i-\textbf{y}oa                  \\
KOL - SW                  & \textbf{j̀}ómì               & \textbf{j}ɛnɛ́                    & \textbf{y}oa                    \\
BUB - SW                  & dì-ómì             & li-ɛɲí                  &                        \\
Kundu - N                 &                       &                          &    \\
   \lspbottomrule
 \end{tabular}
\end{table}

We reconstruct in total 16 PSB consonants: *p, *t, *k, *b, *d, *l, *m, *n, *ɲ, *v, *s, *w, *y, *mb, *nd, *ŋg. The PB reconstructed by \citet{guthriecomparativebantu} and \citet{meeussen1967} are very similar; although they include the addition of PB *g and *c (Guthrie also distinguished *j from *y, see \citealt{wills2022} for details). In these PB system reconstructions, there are no fricative consonants reconstructed (see \citealt{bostoen2019handbook} for more details on PB). 

\largerpage
\citet{hedinger1987} reconstructed for PM in C1 and C2 positions almost the same set of consonants with *p, *t, *k, *s, *b, *d, *j,*m, *n, *ny, *mb, *nd, *ŋg and *w. He also had *f (but not *v) and he reconstructed *g. A closer look at Hedinger’s data showed that the reflexes of PM *g in C1 position are only attested when preceded by a nasal prefix \citep[39]{hedinger1987}. However, in C2, we observe that PM *g is lenited in some languages displaying either /ʔ/, /k ̚ / or /g ̚ /. The choice made by Hedinger to reconstruct *g in PM is understandable as we find attestations of /g/ in C1 and C2 positions as opposed to Sawabantu languages where /g/ is absent in C2 (for example, we have in PM `penis' *jòg ́ > [n-jɔg ̚ ] in Akossi whereas we have *yoko in PSB). 

The comparison of the two proto-systems reconstructed for Manenguba and Sawabantu (languages spoken in the same region, e.g., \figref{fig:grollemundbrownrule:languagearea} and classified within the same subgroup, e.g., \figref{fig:grollemundbrownrule:comparison}) illustrates some similarities between the two subgroups; however, we also observe some key differences. For example, the same sets of consonants have been reconstructed in both proto-languages (with the exception of *g and *v) but the study of the reflexes within the daughter languages shows different evolution patterns for the same proto-consonants in PM and PSB. The examination of \tabref{tab:grollemundbrownrule:manenguba1} displaying the reflexes of PM’s consonants (without the nasals) in C1 position shows that Manenguba languages seem more conservative (with respect to the protolanguage) than Sawabantu languages where multiple reflexes are attested for the same reconstructed consonants (e.g., PSB *d > /ɗ/, /d/, /l/ or /r/). However, the loss of *V2 in Manenguba languages led to the weakening of *C2 triggering the loss of the voicing distinction for plosives whereas in Sawabantu languages, V2 is maintained (along with the voicing distinction). Finally, vowels in Manenguba languages underwent many changes \citep[44-45]{hedinger1987} as opposed to Sawabantu languages. 

\begin{table}
\caption{Reflexes of PM’s plosives, fricatives and the approximant lateral in C1 position (Hedinger 1989: 38-39).}
\label{tab:grollemundbrownrule:manenguba1}
 \begin{tabular}{llllllllll}
  \lsptoprule
\textbf{PM} & \textbf{*p$_1$} & \textbf{*t$_1$} & \textbf{*k$_1$} & \textbf{*b$_1$} & \textbf{*d$_1$} & \textbf{*g$_1$} & \textbf{*s$_1$} & \textbf{*f$_1$} & \textbf{*l$_1$} \\
\midrule
MBM         & p            & t            & k            & b            & d            & g            & s            & f / h        & l / d        \\
MBM         & p            & t            & k            & b            & d            & g            & s / ʃ        & f            & l / d        \\
MYE         & p            & t            & k            & b / bɦ       & d            & g            & s            & h            & l            \\
MBE         & p            & t            & k            & b            & d            & g            & s / ʃ        & h            & l            \\
ELU         & p            & t            & k            & b            & d            & g            & s            & h            & l            \\
NNE         & p            & t            & k            & b            & d            & g            & s            & h            & l            \\
AKO         & p            & t            & k            & b            & d            & g            & s            & h            & l            \\
MHE         & p            & t            & k            & b            & d            & g            & s            & h            & l            \\
MWK         & p            & t            & k            & b            & d            & g            & s            & h            & l            \\
MKA         & p            & t            & k            & b            & d            & g            & s            & h            & l            \\
BLN         & p            & t            & k            & b            & d            & g            & s            & h            & l            \\
BBO         & p            & t            & k            & b            & d            & g            & s            & h / f / v    & y            \\
LEF         & p            & t            & k            & b            & d            & g            & s / ʃ        & h            & l / d        \\
LEK         & p            & t            & k            & b / p        & d            & g            & s            & h            & l / d            \\ 
   \lspbottomrule
 \end{tabular}
\end{table}

\subsection{Sound changes in Sawabantu languages and preliminary observations} \label{sect:browngrollemundrule:observations}

The sound changes observed in \tabref{tab:grollemundbrownrule:consonants}, \tabref{tab:grollemundbrownrule:voicedplosivesexamples}, \tabref{tab:grollemundbrownrule:fricativesc1andc2}, \tabref{tab:grollemundbrownrule:nasalscorrespondence}, and \tabref{tab:grollemundbrownrule:approximants} have shown that Sawabantu languages underwent many changes with several PSB consonants undergoing lenition. The analysis of sound changes supports three genetic subgroups in accordance with the three geographical subgroups distinguished in \citet{ebobisse1989} and \citet{ebobisse2015} with the addition of Kundu (North) with (i) South: PSB *b$_{12}$ > /ɓ/, PSB *d$_{12}$ > /ɗ/ and PSB *s$_{12}$ > /h/, (ii) Central: we have a phonological conditioning found for PSB *b$_{12}$ and PSB *d$_2$ with also PSB *k$_{12}$ > Ø in Duala, Bidiman, Oli, Pongo and Mulimba, and (iii) Southwest and North: PSB *d$_{12}$ > /l/, the lenition of PSB *p$_{12}$ > /ɸ/ and PSB *b$_{12}$ > /β/. Therefore, the definition of our sub-groups differs from Ebobissé’s (\citeyear{ebobisse1989}) analysis in so far as he did not include Oroko (aka North) and so did not notice the proximity between the latter and Southwest.

In addition, the treatment of PSB *d$_{12}$ (and *l$_{12}$) is paramount in Sawabantu languages for possibly establishing some putative chronology. With the three genetic subgroups distinguished (South, Central and Southwest+North), how PSB *d$_{12}$ evolved in these languages is informative. For Southwest languages, PSB *d$_{12}$ and *l$_{12}$ have merged into /l/. This phonological merger defining Southwest and North languages must have occurred last because otherwise, all the PSB *d would be realized /l/ in the remaining subgroups. Instead, the distinction between PSB *d$_{12}$ and *l$_{12}$ is maintained in Central and South languages. The retention of PSB *d$_1$ in Central languages along with the conditioning observed in some Central languages for PSB *d$_2$ as opposed to the unconditioned sound change PSB *d$_{12}$ > /ɗ/ in South languages confirms that we are dealing with two different subgroups. 

Finally, the results observed are in accordance with the groupings distinguished in lexicon-based phylogenies with South and Central languages classified under a same node whereas Southwest and North are classified in a different subgroup linked to Manenguba languages (e.g., \figref{fig:grollemundbrownrule:comparison}).

\subsection{Comparing Proto-Bantu and Proto-Sawabantu: new perspectives} \label{sect:browngrollemundrule:newperspectives}

In this section, we compared our reconstructions established for PSB with PB reconstructions in order to establish sound correspondences between the two proto-languages with a focus on plosive consonants in C1 and C2 positions. The aim is to bring a new perspective on the phenomena described in the previous section by studying Sawabantu languages in relation to Proto-Bantu. 

The examination of \tabref{tab:grollemundbrownrule:manenguba2} shows that PB *p yielded to two reflexes in PSB with *p and *v. \cite[25-26]{philippson2022} mentioned a possible partial conditioning *pVp that could trigger the strong reflexes of PB *p for NW Bantu languages which prevented him from positing double reflexes for PB *p. For PB *t, we have a case of double reflexes or MUR where PB *t > PSB *t or *l with no apparent conditioning. However, \citet[30]{philippson2022} did notice that the strong reflexes of PB *t tend to appear before high vowels in several NW Bantu languages. Our data shows that the reflexes of PSB *t and *l occur in the same phonological contexts confirming the status of double reflexes for PB *t with however one observation: the sequences /liC2V2/, /luC2V2/ or /C1V1li/, /C1V1lu/ are extremely rare in Sawabantu languages. As mentioned in the previous section, PB *k > Ø in Sawabantu languages but also in the remaining NW Bantu languages (A and some B20) which triggered PB *g > *k in the same languages \citep{grollemundandphilippsonforth2}. PB *b is retained in PSB. Finally, PB *d is deleted when followed by non-high vowels and retained (PSB *d) when followed by the vowels *i and *u, a phenomenon observed for many Bantu A languages \citep[30-31]{philippson2022}. 


\begin{table}
\caption{Comparison between PB and PSB for plosives in C1 and C2 positions.}
\label{tab:grollemundbrownrule:manenguba2}
\fittable{ 
 \begin{tabular}{llllllll}
  \lsptoprule
\textbf{PB} & \textbf{PSB}                                                             & \textbf{S}                                                              & \textbf{C}                                                                              & \textbf{SW}                                                                     & \textbf{N}                                                              & \textbf{Examples in PB (from   BLR3)}                                                                                                                                                                                                      & \textbf{PSB}                                                                                                                            \\
\midrule
*p          & \begin{tabular}[c]{@{}l@{}}*p\\    \\ *v\end{tabular}                    & \begin{tabular}[c]{@{}l@{}}p\\    \\ v\end{tabular}                     & \begin{tabular}[c]{@{}l@{}}p\\    \\ w $\sim$v $\sim$f\end{tabular}                     & \begin{tabular}[c]{@{}l@{}}ɸ\\    \\ Ø $\sim$w\end{tabular}                     & \begin{tabular}[c]{@{}l@{}}ɸ\\    \\ Ø\end{tabular}                     & \begin{tabular}[c]{@{}l@{}}PB *pàpá “wing” (2407)\\    \\ PB *pácà “twin” (2348)\\    \\ PB *cʊ́pà “calabash” (746)\end{tabular}                                                                                                      & \begin{tabular}[c]{@{}l@{}}PSB *papa\\    \\ PSB *vasa\\    \\ PSB *sova\end{tabular}                                                   \\
*b          & *b                                                                       & ɓ                                                                       & b                                                                                       & β                                                                               & b $\sim$β                                                               & \begin{tabular}[c]{@{}l@{}}PB *bàdɪ́ “two” (9093)\\    \\ PB *tʊ́ʊ́bá “six” (3034)\end{tabular}                                                                                                                                          & \begin{tabular}[c]{@{}l@{}}PSB *baa\\    \\ PSB *toba\end{tabular}                                                                      \\
*t          & \begin{tabular}[c]{@{}l@{}}*t\\    \\  \\    \\ *l\end{tabular}          & \begin{tabular}[c]{@{}l@{}}t\\    \\  \\    \\ l\end{tabular}           & \begin{tabular}[c]{@{}l@{}}t\\    \\  \\    \\ l\end{tabular}                           & \begin{tabular}[c]{@{}l@{}}t\\    \\  \\    \\ l $\sim$r\end{tabular}           & \begin{tabular}[c]{@{}l@{}}t\\    \\  \\    \\ r\end{tabular}           & \begin{tabular}[c]{@{}l@{}}PB *tʊ́ɪ̀ “ear” (3030)\\    \\ PB *dùt “pull” (1267)\\    \\ PB *téndé “palm tree” (2849)\\    \\ PB *tátʊ̀ “three” (2811)\end{tabular}                                                                     & \begin{tabular}[c]{@{}l@{}}PSB *toi\\    \\ PSB *dut\\    \\ PSB *lɛndɛ\\    \\ PSB *lalo\end{tabular}                                  \\
*d          & \begin{tabular}[c]{@{}l@{}}Ø\\    \\  \\    \\  \\    \\ *d\end{tabular} & \begin{tabular}[c]{@{}l@{}}Ø\\    \\  \\    \\  \\    \\ ɗ\end{tabular} & \begin{tabular}[c]{@{}l@{}}Ø\\    \\  \\    \\  \\    \\ d $\sim$l $\sim$ɗ\end{tabular} & \begin{tabular}[c]{@{}l@{}}Ø\\    \\  \\    \\  \\    \\ l $\sim$r\end{tabular} & \begin{tabular}[c]{@{}l@{}}Ø\\    \\  \\    \\  \\    \\ r\end{tabular} & \begin{tabular}[c]{@{}l@{}}PB *dʊ́mè “husband” (1183)\\    \\ PB *dànd “to buy” (849)\\    \\ PB *búdà “rain” (368)\\    \\ PB *dìbà “water” (1025)\\    \\ PB *dùt “pull” (1267)\\    \\ PB *kúdù “tortoise” (2110)\end{tabular} & \begin{tabular}[c]{@{}l@{}}PSB *ome\\    \\ PSB *and\\    \\ PSB *bua\\    \\ PSB *diba\\    \\ PSB *dut\\    \\ PSB *kudu\end{tabular} \\
*k          & Ø                                                                        & Ø                                                                       & Ø                                                                                       & Ø                                                                               & Ø                                                                       & \begin{tabular}[c]{@{}l@{}}PB *kádí “wife” (1674)\\    \\ PB *Nkákà “pangolin” (1684)\end{tabular}                                                                                                                                     & \begin{tabular}[c]{@{}l@{}}PSB *adi\\    \\ PSB *ka\end{tabular}                                                                        \\
*g          & *k                                                                       & h $\sim$k                                                               & Ø $\sim$k                                                                               & Ø $\sim$k                                                                       & k                                                                       & \begin{tabular}[c]{@{}l@{}}PB *gʊ̀bò “skin” (1479)\\    \\ PB *gòngá “spear” (1448)\\    \\ PB *jògù “elephant” (1607)\\    \\ PB *jɪ́g “learn” (3338)\end{tabular}                                                                   & \begin{tabular}[c]{@{}l@{}}PSB *kobo\\    \\ PSB *kɔngɔ\\    \\ PSB *yɔku\\    \\ PSB *yok\end{tabular}                                     \\ 
   \lspbottomrule
 \end{tabular}
 }
\end{table}


The comparison between PB and PSB shed a new light on the origin of the phoneme /l/ in Sawabantu languages. For Southwest languages, the reflexes /l/ originated from two different PB consonants: *t and *d (reconstructed *l and *d in PSB) which confirms why we had two distinct sound sets in \sectref{sect:browngrollemundrule:consonants}. It also explains why the reflexes of PSB *d can only be found before high vowels as PB *d > Ø before non-high vowels. According to \citet[34]{philippson2022}, the deletion of PB *d before non-high vowels occurred with an intermediate stage attested in some NW Bantu languages with PB *d > /l/ > Ø (whereas PB *d > /d/ or /ɗ/ before high vowels). We found this intermediate stage in some Central Sawabantu languages where we have observed for C2 PSB *d > /l/ before non-high vowels confirming this lenition chain which implied that PB *d$_2$ > /l/ or /ɗ/ or Ø before non-high vowels. 

In addition, according to \citet[37]{philippson2022}, the sound change PB *d > (*l) > Ø must have occurred before PB *t > PSB *l otherwise, most of the /l/ reflexes of PSB *l would have been deleted. Therefore, it seems that the loss of PB *d (> weak reflexes) triggered PB *t > PSB *l (as well as in some other NW Bantu languages, see \citealt[30-31]{philippson2022}) with the result of reintroducing the missing sound /l/. Interestingly, in Sawabantu languages, the reflexes of PSB *l are found in front of non-high vowels which could corroborate that the sound change from PB *t > PSB *l occurred to fill the void created by the rule PB *d > *l > Ø in the same context (but we need to recall that PB *t > PSB *t in all contexts so we do not have an explanation to why some words with PB *t > PSB *l before non-high vowels). Moreover, we can note that some languages such as Mokpe or Bubia present PSB *l > Ø or /l/ which could also indicate that the process of deleting PSB *l is continuing. 

The same phenomenon is observed with PB *k. When PB *k is deleted, Sawabantu languages (as well as the remaining NW Bantu languages A and some B20) reintroduced the missing sound with the sound change PB *g > PSB *k. Interestingly, PSB *k is undergoing lenition again in Sawabantu languages. 

In conclusion, the examination of \tabref{tab:grollemundbrownrule:manenguba2} highlights interesting sound shifts from PB to Sawabantu languages with strategies of replacement for sounds that have been deleted.

\section{Conclusion} \label{sect:browngrollemundrule:conclusion}

Sawabantu languages are fascinating for several reasons highlighted in this study. First, the study and reconstruction of their consonants evidences some interesting phenomena occurring at the PSB level with some PSB consonants undergoing lenition (PSB *p > /ɸ/, *k > Ø, *b > /β/ or /w/, *d > /l/). The study of PSB sound changes divided the Sawabantu languages into three subgroups: South, Central and Southwest-North. At the PB level, some interesting sound changes are observed with PB *t > PSB *t and *l, PB *k > Ø, PB *g > PSB *k and PB *d > Ø when followed by non-high vowels confirming that languages spoken in this particular region underwent many sound changes that resulted in the modification of the shape of their lexicon in comparison with PB (e.g., PB *kútà BLR3 2138 `oil' > PSB *ula but [ma-kuta] in F31 Nilamba (data from \citealt{yukawa1989}) - a Bantu language spoken in Tanzania, PB *dʊ́k BLR3 1179 “vomit” > PSB *o but [kwi-lʊka] in F31 Nilamba). Finally, this study confirms the different status of Manenguba languages (A15) that underwent different sound changes in C1 and C2 positions which led \citet{nurseandphilippson2003} to exclude these languages from the Sawabantu clade. 

\section*{Acknowledgements}

We wish to thank the two anonymous reviewers for their useful feedback. First author’s research was supported by the National Science Foundation (NSF), award 2152822.

\sloppy %this command eliminates a quirk that appears in the bibliography, you can just leave it here.

\printbibliography[heading=subbibliography,notkeyword=this]

\end{document}
